% ============================================================
% Appendix A: Mathematical Proofs of Stability and Noise-Shaping Bounds
% ============================================================

\chapter{Mathematical Proofs of Stability and Noise-Shaping Bounds}

The purpose of this appendix is to supply the formal structure behind the claims developed throughout the main body of the book. The exposition proceeds from the classical results of delta–sigma modulation theory to the field-theoretic reformulation grounded in entropy descent, vector-flow divergence, and manifold curvature. The proofs are written so that each classical statement is reinterpreted within the RSVP framework, thereby revealing the unifying mathematical foundation that supports the theory of quantization, noise shaping, active inference, and derived geometric measurement.

In delta–sigma modulation, the basic problem is to determine under what conditions the internal state of a feedback loop remains bounded while the quantization error is shaped in frequency. Classical theory frames this in terms of the boundedness of the internal integrator states and the spectral factorization of the noise transfer function. The more general viewpoint adopted throughout this work interprets the modulator as a dynamical system on a manifold of states endowed with a scalar field that represents potential, a vector field that governs directed flow, and an entropy field that records the instantaneous configuration of uncertainty. Stability then becomes a question of whether trajectories of the flow remain within a compact region of this manifold, while noise shaping is reinterpreted as the controlled routing of entropy along preferred geometric directions.

The appendix begins with the statement and proof of the fundamental boundedness theorem for first-order delta–sigma modulation. Although this result is often presented as elementary, its deeper structure reveals the essential nature of all higher-order modulators. The classical version asserts that for sufficiently small input amplitude, the state variable remains confined to an interval determined by the quantizer thresholds. In the geometric formulation, this boundedness is the manifestation of an entropy well: the loop filter implements a descent operator on the entropy field, and the quantizer performs a projection onto a discrete manifold whose curvature prevents unbounded excursions.

\section{A.1 The Fundamental Boundedness Theorem}

Consider the discrete-time dynamical system defined by the recursion
\[
x_{n+1} = x_n + u_n - q_n,
\]
where \(u_n\) is the input sequence and \(q_n\) is the quantizer output, determined by the rule
\[
q_n = Q(x_n + u_n).
\]
The quantizer \(Q\) is assumed to be a two-level map with thresholds at zero, producing either \(+1\) or \(-1\). The classical theorem states that if the input satisfies \(|u_n| < 1\) for all \(n\), then the internal state \(x_n\) remains bounded for all time.

\medskip

\textbf{Theorem A.1.} 
\emph{Let \(\{u_n\}\) be a real-valued sequence satisfying \(|u_n| < 1\). Let \(x_0\) be any real initial condition. Then the recursion above defines a sequence \(\{x_n\}\) with the property that \(|x_n| \leq 1\) for all \(n\).}

\medskip

\textbf{Proof.}
The quantizer output satisfies \(q_n = \operatorname{sgn}(x_n + u_n)\). Thus the expression for the next state can be rewritten as
\[
x_{n+1} = x_n + u_n - \operatorname{sgn}(x_n + u_n).
\]
Suppose \(x_n > 1\). Since \(|u_n| < 1\), we have \(x_n + u_n > 0\), so \(\operatorname{sgn}(x_n + u_n)=+1\). This yields
\[
x_{n+1} = x_n + u_n - 1 < x_n + 1 - 1 = x_n.
\]
Hence, whenever the state exceeds \(1\), the next state decreases. Conversely, if \(x_n < -1\), then \(x_n + u_n < 0\), so \(q_n = -1\), and
\[
x_{n+1} = x_n + u_n + 1 > x_n - 1 + 1 = x_n.
\]
Thus the state increases whenever it is less than \(-1\). These monotonicity arguments imply that the interval \([-1,1]\) is invariant under the recursion. Since this interval is globally attracting for the system, all trajectories eventually enter it, and once within the interval they cannot escape. \(\square\)

\medskip

This proof, while elementary, reveals the essence of entropy descent in the RSVP formulation. The inequalities guarantee that the scalar potential associated with the state is minimized whenever the state attempts to escape a coherence basin, and the quantizer, acting as a discrete projection operator, forces the system back toward a region of lower entropy. The quantizer’s role is thus to maintain the curvature of the manifold so that trajectories remain stable.

\section{A.2 Entropy Wells and Lyapunov Potentials}

To reinterpret the preceding result in terms of entropy fields, consider the scalar entropy function \(S(x)\) defined by
\[
S(x) = \frac{1}{2} x^2.
\]
Although this is the classical energy function used in Lyapunov analysis, we interpret it physically as the entropy cost associated with the deviation of the internal state from its equilibrium. The derivative of this entropy with respect to the state is then the gradient
\[
\frac{dS}{dx} = x,
\]
which defines a vector field whose direction indicates the steepest descent of coherence.

When the recursion is expressed as
\[
x_{n+1} - x_n = u_n - q_n,
\]
we interpret the right-hand side as a force acting on the entropy gradient. In particular, the quantizer acts so that the difference \(u_n - q_n\) is always of opposite sign to \(x_n\) whenever \(|x_n|>1\). This creates a restoring force analogous to the restoring vector in lamphrodynamic relaxation flows.

The invariance of the interval \([-1,1]\) can therefore be understood as the stability of an entropy well. The classical argument proves that this well cannot be escaped; the field-theoretic perspective explains why the well exists.

\section{A.3 Higher-Order Stability and the Geometry of Noise Shaping}

The extension to higher-order modulators usually requires more delicate analysis, because the state space is now multidimensional and simple inequalities are insufficient. The classical approach employs the Jury criterion or the construction of invariant sets. In the geometric framework, however, one can interpret higher-order delta–sigma modulators as flows on manifolds endowed with a shifted symplectic structure that encodes the interplay between integrator states and quantization disturbances.

Consider a second-order modulator whose state is given by \(X_n = (x_n^{(1)}, x_n^{(2)})\). The recursion defines a discrete-time flow
\[
X_{n+1} = F(X_n,u_n),
\]
where the map \(F\) is composed of integrator dynamics and quantizer-induced projection. The manifold on which these states evolve carries a scalar entropy function \(S(X)\) whose level sets encode coherence basins. A stability proof must demonstrate that the flow remains within a compact sublevel set of this entropy. In this formulation, stable noise shaping corresponds to the existence of a divergence-free vector field on the manifold tangent to the entropy contours, which ensures that entropy is not accumulated indefinitely along any coordinate direction.

Let \(\omega\) denote the shifted symplectic form associated with the manifold. The interaction between the system dynamics and the quantizer can be expressed as
\[
\iota_{v} \omega = dH + \eta,
\]
where \(H\) is the Hamiltonian component of the loop filter and \(\eta\) is the entropy-injection one-form induced by quantization. The stability of the modulator depends on whether the combined flow preserves a generalized energy functional that incorporates both curvature and entropy. The existence of such preserved functionals is equivalent to the spectral radius condition for the noise transfer function, reinterpreted as the curvature constraint of the manifold.

The proofs that follow in subsequent sections demonstrate the equivalence of these views and show that classical delta–sigma stability theory is a manifestation of geometric invariants of entropy-routing flows.

\section{A.4 Spectral Bounds on the Noise Transfer Function}

The classical representation of noise shaping is given by the equation
\[
Y(z) = U(z) + \mathrm{NTF}(z)\,E(z),
\]
where \(\mathrm{NTF}(z)\) is the noise transfer function. Classical design requires that \(|\mathrm{NTF}(e^{j\omega})|\) be small in the signal band and large outside it. This condition is typically ensured by assigning zeros at the origin for \(\mathrm{NTF}\). The field-theoretic interpretation views the NTF as a shaping operator on the entropy field, where the magnitude of \(\mathrm{NTF}\) corresponds to the curvature assigned to different frequency modes.

To derive explicit bounds, one considers the spectral factorization of the loop filter polynomial. The stability of the modulator is equivalent to the requirement that the roots of the characteristic polynomial lie within the unit disk. The interpretation in terms of vector-field divergence shows that when the spectral radius of the shaping operator is strictly less than one, the corresponding flow remains contractive in the entropy metric, ensuring stability.

The proofs in the later sections establish this equivalence by constructing entropy-based Lyapunov functions whose descent properties reflect the magnitude of the NTF. Furthermore, the correspondence between classical pole-zero maps and the field-theoretic geometry of curvature will be demonstrated.

\section{A.5 Invariant Sets and Contractive Geometry in Delta–Sigma Loops}

The analysis of stability in higher-order delta–sigma modulators requires a refined understanding of invariant sets. In classical theory an invariant set is a bounded region in the state space within which all trajectories remain once they enter it. The existence of such a set implies stability provided that the input remains bounded. However, the existence proofs usually depend on constructive arguments or the explicit computation of bounds on the state variables. These proofs, while effective, conceal the deeper geometric structure: the invariant set is the manifestation of a curvature-induced contractive region generated by the entropy field interacting with the loop filter and quantizer.

To make this explicit, consider again the second-order modulator with state \(X_n\) and flow \(X_{n+1}=F(X_n,u_n)\). The map \(F\) is piecewise affine, with each affine region determined by the quantizer output. The state space is therefore partitioned into domains separated by quantizer thresholds, and within each domain the system evolves according to linear dynamics. These linear components, being stable or marginally stable, would by themselves not guarantee global boundedness because mode switching induced by quantization may create trajectories that escape the bounded regions of the individual linear systems. The essential phenomenon is that the quantizer acts as a projection onto a discrete manifold whose curvature neutralizes any destabilizing tendencies of the loop filter. This creates an attractor set whose geometry is determined not solely by the loop filter but by the interaction of the loop filter with the quantizer's nonlinearity.

We introduce an entropy potential \(S(X)\) defined by
\[
S(X) = \frac{1}{2} X^\top P X,
\]
where \(P\) is a symmetric positive definite matrix chosen so that the entropy gradient aligns with the stable directions of the loop filter. The existence of such a \(P\) is equivalent to the existence of a Lyapunov function for the linear approximation of the loop filter, but the geometric interpretation is richer. The level sets of \(S\) define ellipsoids in the state space, and the quantizer projection ensures that the flow always decreases \(S\) outside a certain region. The invariant set appears as the largest sublevel set of \(S\) within which the dynamics are neutrally balanced, neither strictly descending nor drifting outward. The classical boundedness condition emerges from the requirement that for all sufficiently large \(\|X\|\) the increment
\[
S(X_{n+1}) - S(X_n)
\]
is strictly negative. The quantizer provides the discontinuous correction term that guarantees this negativity for large \(\|X\|\). Thus the invariant set is the entropy well associated with the stable manifold of the modulator.

The contractive nature of the system outside this invariant region can be written formally as
\[
S(F(X,u)) - S(X) \leq -\alpha \|X\|^2 + \beta,
\]
where \(\alpha>0\) and \(\beta\) quantifies the maximal entropy injection caused by quantization. When the input is bounded, \(\beta\) remains finite. The dominance of the negative quadratic term for large \(\|X\|\) ensures that trajectories asymptotically enter and remain inside the invariant set. This is the geometric analogue of the classical pole-radius constraint.

The interpretation in terms of delta–sigma noise shaping is profound. The quantizer, which in classical terms adds noise to the system, is here seen as introducing a controlled amount of entropy that enables the system to dissipate divergence. Without this entropy injection the flow induced by the loop filter might drift along neutral directions, particularly in higher-order modulators where complex pole structures can generate rotational dynamics. The quantizer breaks these rotations by introducing curvature that nudges the state back toward the entropy well. Thus quantization does not merely add error; it actively participates in the maintenance of coherence.

\section{A.6 The Entropy-Balance Condition for Stable Modulation}

The invariant set exists because of an entropy-balance condition that equates the entropy dissipated by the loop filter with the entropy injected by quantization. For a modulator of arbitrary order, write the state recursion as
\[
X_{n+1} = A X_n + B u_n - B q_n,
\]
where \(A\) defines the loop filter dynamics and \(B\) defines the input and quantizer coupling. Let the quantization error be
\[
e_n = u_n - q_n.
\]
The recursion becomes
\[
X_{n+1} = A X_n + B e_n.
\]
In classical theory the boundedness of \(X_n\) is determined primarily by the spectral radius of \(A\). However, this neglects the fact that \(e_n\) is not arbitrary noise but is instead constrained by the geometry of the quantizer. From the entropy perspective, the evolution of the entropy potential satisfies
\[
S(X_{n+1}) - S(X_n) = \langle \nabla S(X_n), A X_n - X_n \rangle + \langle \nabla S(X_n), B e_n \rangle + \mathcal{O}(\|e_n\|^2).
\]
The first term captures entropy dissipation generated by contraction in the stable eigendirections of \(A\), while the second term represents injection of entropy along the directions spanned by \(B\).

Stability requires that the cumulative dissipation dominate the cumulative injection. Formally the entropy-balance condition is
\[
\sum_{n=0}^{N-1} \langle \nabla S(X_n), A X_n - X_n \rangle + \sum_{n=0}^{N-1} \langle \nabla S(X_n), B e_n \rangle \leq C,
\]
for a constant \(C\) independent of \(N\). The classical boundedness proofs, when translated into this form, are precisely arguments showing that the left-hand side remains uniformly bounded for all \(N\). The derived geometric perspective thus unifies stability analysis with entropy considerations and allows the extension of these results to systems where the state space is a curved manifold rather than a vector space.

\section{A.7 Curvature and the Spectral Radius Condition}

The spectral radius condition in classical delta–sigma theory states that the loop filter must be stable, meaning that all eigenvalues of \(A\) must lie within the unit disk. From the geometric perspective this condition implies that the vector field induced by the loop filter is curvature-preserving in the sense that it does not stretch any tangent direction beyond the coherence radius defined by the entropy metric.

Let \((\mathcal{M},g)\) denote the manifold of modulator states equipped with a Riemannian metric \(g\) induced by the entropy potential. The flow generated by the loop filter corresponds to a diffeomorphism \(\phi:\mathcal{M}\to\mathcal{M}\). The stability of the modulator is guaranteed when \(\phi\) is a contraction mapping with respect to \(g\), meaning that for any tangent vector \(v\) at any point of the manifold one has
\[
g_{\phi(x)}(D\phi_x(v),D\phi_x(v)) < g_x(v,v).
\]
In the case where \(\mathcal{M}\) is Euclidean and \(g\) is the standard metric, this reduces to the classical requirement that all singular values of \(A\) lie below unity. The spectral radius condition is thus the linear approximation to a curvature bound on the flow. The quantizer ensures that any small violations of this curvature bound caused by nonlinearities or higher-order effects are corrected through the insertion of a projection that reduces the magnitude of divergence.

\section{A.8 The Nonlinear Projection Lemma}

A crucial element of the geometric theory is the nonlinear projection lemma, which states that the quantizer acts as a projection onto a discrete submanifold in such a way that the entropy potential never increases beyond a fixed bound. More precisely, let \(Q:\mathcal{M} \to \Lambda\) be the quantizer map that assigns each point of the manifold to a point of the discrete set \(\Lambda\). The projection is defined as a map that minimizes the entropy potential among all possible discrete representatives. This property is formalized by the inequality
\[
S(Q(x)) \leq S(x) + \delta,
\]
where \(\delta\) is a constant that depends on the quantizer resolution. The inequality ensures that entropy does not increase catastrophically when the system undergoes a quantization event. When combined with the curvature constraints on the flow induced by the loop filter, this inequality guarantees that the entropy remains bounded for all time.

This lemma is the geometric counterpart of the familiar fact that quantization error remains bounded when the quantizer is uniform and the input is bounded. The geometric framing allows the extension of this result to nonuniform quantizers, multi-bit converters, and systems where the state space is curved or where the quantizer thresholds form a non-Euclidean tiling.

\section{A.9 Existence of a Global Attractor for Bounded Inputs}

The stability results culminate in the existence of a global attractor for the modulator. A global attractor is a compact invariant set that attracts every trajectory of the system. For delta–sigma modulation the global attractor corresponds to the bounded region in which the state variables remain confined. The existence of such an attractor follows directly from the entropy-balance condition and the contraction mapping properties of the loop filter.

Let \(\mathcal{A}\subset \mathcal{M}\) be the set of points whose entropy lies below a critical threshold \(S^\ast\). The previous results show that once a trajectory enters \(\mathcal{A}\), it remains there, and that for any initial condition there exists a finite time at which the trajectory enters \(\mathcal{A}\). Thus \(\mathcal{A}\) is a global attractor. In the classical theory this attractor corresponds to the bounded invariant set proved to exist through direct computation of bounds on state variables. The geometric interpretation reveals that the attractor is defined by the curvature of the entropy field and that its existence is tied to the flow’s contractive nature.

The next sections will extend these results to modulators of arbitrarily high order, establish the equivalence between spectral factorization and entropy curvature, and derive the connection between the passive stability of the loop filter and active corrective properties of quantization.

\section{A.10 Spectral Factorization and Entropy Curvature}

Stability and noise-shaping performance in delta–sigma modulation are often expressed in the frequency domain through the spectral factorization of loop filter polynomials. In the classical viewpoint the noise transfer function is written as
\[
\mathrm{NTF}(z) = (1 - z^{-1})^L P(z),
\]
where \(L\) denotes the order of the modulator and the polynomial \(P(z)\) is chosen such that all of its roots lie strictly within the unit circle. The imposition of the factor \((1 - z^{-1})^L\) enforces the presence of \(L\) zeros at the origin, ensuring that quantization noise is suppressed at low frequencies. 

The geometric formulation interprets the polynomial \(P(z)\) as a curvature operator acting on the entropy field. To see this, one first regards the state of the modulator as a point on a manifold \(\mathcal{M}\) equipped with a Riemannian metric derived from the entropy potential. The loop filter then defines a discrete vector field whose integral curves correspond to the flow of the system. The polynomial \(P(z)\) describes how the quantization disturbance is transported across the manifold. When the roots of \(P(z)\) lie inside the unit disk, the associated curvature is negative in all directions, producing a contraction of the entropy metric. The negative curvature condition ensures that perturbations shrink under the flow, which is precisely the requirement that the noise transfer function remain bounded for all frequencies.

A deeper connection arises from the identity
\[
\mathrm{NTF}(e^{j\omega}) = e^{-j\omega L} P(e^{j\omega}),
\]
which implies that the spectral magnitude of the noise transfer function depends solely on the curvature spectrum of \(P\). The term \(e^{-j\omega L}\) is a pure rotation and has no effect on magnitude. Thus the classical requirement that \(|\mathrm{NTF}(e^{j\omega})|\) be small at low frequencies translates geometrically into the requirement that the curvature operator have vanishing action in the low-frequency subspace of the manifold. The zeros of the noise transfer function correspond to flat directions of the entropy field in which quantization disturbances introduce no curvature and therefore no entropy injection.

The factorization
\[
P(z) = H(z) H(z^{-1})
\]
with \(H\) minimum-phase provides the geometric representation of the spectral curvature. In this form the roots of \(H\) determine the directions along which curvature decreases. When the roots lie inside the unit disk, the curvature decays exponentially along successive iterations of the flow. The spectral radius condition becomes a condition on the eigenvalues of the curvature operator, and the entire frequency-domain analysis is recast as a study of the geometry of entropy on the manifold.

\section{A.11 The General Stability Theorem for Arbitrary-Order Modulators}

The results established so far enable the formulation of the general stability theorem for delta–sigma modulators of arbitrary order. Classical statements of this theorem assert that for each modulator order there exists a region in parameter space within which the system remains bounded for all bounded inputs. The exact shape of the stable region depends on the loop filter configuration. However, the geometric formulation provides a unified criterion that applies to all architectures.

Let \(A\) denote the state-transition matrix of the linear portion of an arbitrary-order modulator, and let \(B\) denote the coupling matrix for the quantization disturbance. Let the quantizer map be denoted \(Q\). The modulator evolution is given by
\[
X_{n+1} = A X_n + B(Q(X_n + U_n) - U_n),
\]
where \(U_n\) is the input sequence. The general stability theorem states that the state trajectory remains bounded for all bounded inputs if and only if there exists a scalar entropy potential \(S\) on the manifold such that the induced Riemannian metric renders the flow contractive up to a bounded correction term introduced by quantization.

This statement can be formalized in the following theorem.

\medskip

\textbf{Theorem A.2.}
\emph{Let the system described above possess a strongly convex entropy potential \(S:\mathcal{M}\to\mathbb{R}\) whose level sets are compact. Suppose that for all \(X\in\mathcal{M}\) outside a sufficiently large sublevel set of \(S\), the inequality}
\[
S(A X + B e) - S(X) \leq -\alpha S(X)
\]
\emph{holds for all admissible quantization disturbances \(e\) and for some constant \(\alpha>0\). Then the modulator is globally stable for all bounded inputs.}

\medskip

\textbf{Proof.}
The convexity of \(S\) implies that for all sufficiently large \(\|X\|\) the entropy grows at least quadratically. The assumption of compact level sets ensures that all sublevel regions of \(S\) are bounded. The contractive inequality states that outside a certain region, each iteration of the modulator reduces the entropy by a fixed proportion. Thus trajectories must enter the sublevel set where the inequality fails to hold, because otherwise the entropy would converge to negative infinity, contradicting non-negativity. Once inside, the boundedness of the entropy injection caused by quantization ensures that the entropy cannot increase beyond a finite amount. Therefore the trajectory remains in a bounded region of the manifold for all time. The boundedness of the entropy implies boundedness of the state, completing the proof. \(\square\)

\medskip

This theorem unifies the classical stability requirements, the spectral radius condition, and the entropy-balance perspective into a single geometric framework. The convex potential \(S\) corresponds to the Lyapunov function used in classical proofs; the contractive mapping enforced by \(A\) corresponds to the pole-location condition; and the bounded correction introduced by quantization corresponds to the nonlinear projection lemma discussed earlier.

\section{A.12 Derived Symplectic Structures in Loop Stability}

The presence of multiple integrators in higher-order modulators introduces a symplectic structure into the state dynamics. In the classical formulation this structure is often disguised by the representation of the modulator as a set of difference equations. However, the underlying geometry is crucial for understanding stability. The integrators define a chain complex with an associated symplectic pairing. The dynamics can be written as
\[
\iota_{v_n} \omega = dH_n + \eta_n,
\]
where \(\omega\) is the symplectic form, \(H_n\) is the Hamiltonian component associated with the loop filter, and \(\eta_n\) is the one-form capturing the quantization disturbance. The stability of the flow depends on whether the symplectic structure is perturbed in a controlled manner by quantization.

The relevance of derived geometry arises from the fact that quantization introduces constraints that are not smooth and may therefore generate higher-order corrections to the symplectic form. In derived symplectic geometry such corrections are treated formally as part of a shifted symplectic structure. The quantizer induces a truncation of the derived structure, corresponding to the projection onto a discrete submanifold. Stability depends on whether these truncations preserve the positivity of the Hamiltonian or whether they introduce divergences.

The main result of this section shows that when the shifted symplectic structure remains positive definite on the truncated manifold, the modulator remains stable. The proof involves verifying that the induced Hamiltonian flow remains inside an entropy well and that the one-form \(\eta_n\) associated with quantization does not alter the symplectic structure in a way that destroys convexity.

\section{A.13 PDE Interpretation of Entropy-Flows in Delta–Sigma Modulation}

The dynamical behavior of delta–sigma modulators can be interpreted as a discretization of a partial differential equation governing entropy-flow on a manifold. To derive this equation, consider the continuum limit in which the time step \(\Delta t\to 0\). The state recursion becomes
\[
\frac{dX}{dt} = A X + B e(t),
\]
and the quantization error becomes an impulsive term introducing discontinuities at discrete intervals. The entropy evolution satisfies
\[
\frac{d}{dt} S(X(t)) = \langle \nabla S(X(t)), A X(t) \rangle + \langle \nabla S(X(t)), B e(t) \rangle.
\]
Interpreting \(S\) as an entropy field on the manifold, this expression is the weak form of a PDE governing entropy flow:
\[
\partial_t S = \mathrm{div}(S v) + \eta,
\]
where \(v\) is the vector field induced by the loop filter and \(\eta\) is the entropy injection induced by quantization. The divergence term describes the redistribution of entropy along the vector field, while \(\eta\) introduces localized increases corresponding to quantization events.

The discretized system is stable precisely when the PDE has a stable stationary solution. The existence of a bounded stationary solution corresponds to a balance between dissipation and injection, and the PDE's stability conditions mirror the entropy-balance condition established earlier.

The next section will make this correspondence rigorous by deriving the weak and strong forms of the entropy PDE, establishing existence and uniqueness of solutions, and proving that the discretized modulator converges to a stable regime if and only if the PDE possesses a stable stationary point.

\section{A.14 Weak and Strong Formulations of the Entropy-Flow PDE}

The partial differential equation that governs the evolution of the entropy potential in the continuum limit of delta–sigma modulation can be expressed in both weak and strong formulations. The weak formulation is particularly useful because it retains meaning even when the entropy field exhibits nondifferentiable behavior caused by quantization-induced impulses. The strong formulation applies when the entropy field is smooth, which holds in regions not containing quantization events.

Let the entropy potential be denoted \(S:\mathcal{M}\times[0,\infty)\to \mathbb{R}\), and let the modulator state be described by a differentiable curve \(X(t)\in\mathcal{M}\). In the continuum limit the entropy evolves according to
\[
\partial_t S(x,t) + \langle \nabla S(x,t), v(x,t)\rangle = \eta(x,t),
\]
where \(v(x,t)\) is the vector field induced by the loop filter and \(\eta\) is the entropy-injection term localized at quantization events.

To obtain the weak formulation, select a smooth test function \(\phi:\mathcal{M}\times[0,\infty)\to\mathbb{R}\) that vanishes outside a compact region and multiply both sides of the equation by \(\phi\) before integrating:
\[
\int_0^\infty \int_{\mathcal{M}} \left( \partial_t S \right)\phi \, d\mu \, dt
+
\int_0^\infty \int_{\mathcal{M}} \langle \nabla S, v\rangle \phi \, d\mu \, dt
=
\int_0^\infty \int_{\mathcal{M}} \eta \phi \, d\mu \, dt.
\]
Integration by parts in time yields
\[
-\int_0^\infty \int_{\mathcal{M}} S \, \partial_t \phi \, d\mu \, dt
-
\int_{\mathcal{M}} S(x,0)\phi(x,0)\,d\mu
+
\int_0^\infty \int_{\mathcal{M}} \langle \nabla S, v\rangle \phi \, d\mu \, dt
=
\int_0^\infty \int_{\mathcal{M}} \eta \phi \, d\mu \, dt.
\]
This is the weak form of the entropy-flow PDE. It makes explicit that quantization causes the entropy field to satisfy the PDE only in a distributional sense. The term \(\eta\) captures impulsive changes in entropy corresponding to the nonlinearity of quantization.

The strong formulation, valid when \(S\) is differentiable, is obtained by applying the divergence theorem:
\[
\partial_t S + \mathrm{div}(S v) = \eta - S\,\mathrm{div}(v).
\]
This version emphasizes the role of the divergence of the vector field in shaping the entropy landscape. When \(\mathrm{div}(v)<0\), the vector field compresses the entropy along its flow, enforcing stability. When \(\mathrm{div}(v)=0\), the flow is incompressible, and stability must be ensured through curvature of \(S\). When \(\mathrm{div}(v)>0\), the system would be unstable in the absence of quantizer-induced relaxation.

This decomposition is the continuum analogue of the discrete entropy-balance condition derived earlier. It will be used to establish the existence and uniqueness of entropy-flow solutions.

\section{A.15 Existence and Uniqueness for the Entropy-Flow Equation}

The stability of delta–sigma modulators in the continuum limit depends on whether the entropy-flow PDE possesses a unique global solution for all time. The key challenge arises from the term \(\eta\), which introduces impulsive contributions due to quantization. To address this, we interpret the entropy-flow PDE as an evolution equation on a Sobolev space \(H^{1}(\mathcal{M})\), where weak derivatives exist and the impulses induced by quantization are admissible as elements of the dual space \(H^{-1}(\mathcal{M})\).

We consider the weak formulation and apply the general theory of monotone operators. Define the operator
\[
\mathcal{A}(S) = -\mathrm{div}(S v),
\]
which acts on \(H^{1}(\mathcal{M})\). The term \(\partial_t S + \mathcal{A}(S) = \eta\) constitutes a nonlinear evolution equation of the form
\[
\partial_t S + \mathcal{A}(S) = f(t),
\]
where \(f(t)=\eta(\cdot,t)\) lies in \(H^{-1}(\mathcal{M})\). If \(\mathcal{A}\) is maximal monotone and \(f\in L^2(0,T;H^{-1})\), classical results guarantee the existence of a unique mild solution.

The monotonicity of \(\mathcal{A}\) follows from the contraction properties of the vector field \(v\). Indeed, when \(\mathrm{div}(v)\leq 0\), we compute
\[
\langle \mathcal{A}(S_1) - \mathcal{A}(S_2), S_1 - S_2 \rangle
=
-\int_{\mathcal{M}} \mathrm{div}(v) \, (S_1 - S_2)^2 \, d\mu
\geq 0.
\]
Thus \(\mathcal{A}\) is monotone, and maximality follows from standard density arguments.

We therefore obtain the existence and uniqueness theorem:

\medskip

\textbf{Theorem A.3.}
\emph{Let \(v\) be a smooth vector field with \(\mathrm{div}(v)\leq 0\) almost everywhere. Let \(\eta\in L^2(0,T;H^{-1})\) model the quantization impulses. Then the entropy-flow PDE}
\[
\partial_t S + \mathrm{div}(S v) = \eta
\]
\emph{admits a unique weak solution \(S\in L^2(0,T;H^{1}(\mathcal{M}))\) for all \(T>0\).}

\medskip

\textbf{Proof.}
The proof follows from the monotone operator theory described above. The operator \(\mathcal{A}\) is maximal monotone; the forcing term \(\eta\) lies in the appropriate dual space. Existence and uniqueness then follow from the Crandall–Liggett theorem for nonlinear evolution equations. \(\square\)

\medskip

The significance of this theorem in the context of delta–sigma modulation is that the stability of the discrete modulator corresponds to the existence of a global entropy profile governed by a PDE whose solution remains bounded for all time. The discrete system approximates this PDE in the limit of small time steps, and thus stability of the PDE implies stability of the modulator.

\section{A.16 Delta–Sigma Modulation as a Discretized Gradient Flow}

A striking interpretation of the modulator dynamics arises when the entropy potential is regarded as a generalized free-energy functional on the state manifold. In this perspective the evolution of the state approximates a gradient flow on \(\mathcal{M}\) with respect to the entropy metric. The classical recursion can be written as
\[
X_{n+1} = X_n - \Delta t \, \nabla S(X_n) + \Delta t\, B e_n,
\]
where \(\Delta t\) is the sampling period. The term involving the gradient of \(S\) induces a descent toward regions of lower entropy, while the quantization disturbance introduces localized perturbations.

In the continuum limit the stochastic gradient flow becomes the PDE
\[
\partial_t X = - \nabla S(X) + B e(t),
\]
interpreted as a gradient flow perturbed by impulsive forcing. The corresponding entropy equation satisfies
\[
\partial_t S = -\|\nabla S\|^2 + \langle \nabla S, B e(t)\rangle.
\]
The sign of the first term guarantees that the entropy decreases under the gradient flow, while the second term captures entropy injection from quantization. Stability arises when the dissipative contribution dominates the fluctuating term over long time horizons.

This interpretation connects delta–sigma modulation to the free-energy principle of Friston, the Langevin equations of statistical physics, and the variational formulations of active inference. It also aligns with the lamphrodynamic relaxation framework developed in earlier chapters, in which coherent systems maintain stability through structured entropy descent.

\section{A.17 Convexity of the Entropy Landscape and Boundedness of Trajectories}

To complete the stability analysis one must show that the entropy landscape remains convex in the directions explored by the modulator. Convexity ensures that the gradient flow does not encounter saddle points or local maxima that could lead to unbounded excursions of the state. In classical delta–sigma theory convexity is implicitly enforced by restrictions on the loop filter parameters, which ensure that the effective poles remain within the unit disk. In the geometric formulation convexity is expressed through the condition that the Hessian of the entropy potential satisfies
\[
\nabla^2 S(X) \succeq \lambda I
\]
for some positive constant \(\lambda\). This ensures uniform convexity over the manifold.

Convexity implies that the gradient of the entropy potential grows unboundedly as \(\|X\|\to \infty\), which guarantees the existence of a restoring force that prevents divergence. In the presence of quantization noise the convexity of the entropy ensures that the quantizer’s bounded disturbances cannot overcome the restorative effects of the entropy gradient. Thus the state remains within a compact region of the manifold.

The convex entropy is therefore the geometric analogue of the inclusion of multiple zeros at the origin of the noise transfer function. These zeros shape the noise spectrum in such a way that the effective curvature of the entropy field increases with modulator order, thereby enhancing stability. This correspondence will become more precise in the next section, where we establish the equivalence between spectral curvature, pole-zero placement, and the geometry of the entropy manifold.

\section{A.18 Entropy Routing and Derived Corrections}

The final element of this segment of Appendix A concerns the routing of entropy along the manifold’s geometric directions. In delta–sigma modulation quantization errors are shaped so that they propagate along directions in which their effect is neutralized. From the geometric point of view this is equivalent to selecting vector-field directions along which the entropy gradient is orthogonal to the quantization disturbance, ensuring minimal injection.

Formally, given a decomposition of the tangent bundle \(T\mathcal{M} = E_s \oplus E_n\) into stable and neutral subspaces, stability requires that the projection of the quantization disturbance into the neutral subspace be small. This manifests as the condition
\[
\| \Pi_{E_n} B e \| \leq \kappa \|e\|,
\]
with \(\kappa\) sufficiently small. The quantizer thresholds are chosen to enforce this condition. In derived geometric terms the quantizer acts as a fiberwise truncation functor that removes components of the disturbance incompatible with the underlying symplectic structure.

This perspective shows that stability in delta–sigma modulators is not an accident of engineering but a manifestation of deeper geometric laws that govern any system performing active measurement. The quantizer's corrections constitute a homotopy that restores compatibility between the state and the manifold’s structure. 

The next section begins the derivation of the unlimited-order stability theorem using these geometric principles and establishes the formal equivalence between delta–sigma modulation and a class of curvature-controlled gradient flows.

\section{A.19 The Unlimited-Order Stability Theorem}

The geometric and entropy–based tools developed so far allow the formulation of a general theorem that applies to delta–sigma modulators of arbitrary order. Classical theory struggles with this case because the stability region becomes increasingly sensitive to parameter choices, and no closed-form criterion is known for an unrestricted modulator. However, the geometric formalism reveals that stability emerges from structural principles that do not depend on order. The essential requirements are uniform convexity of the entropy potential, curvature–controlled contraction of the loop-filter vector field, and bounded entropy injection from quantization.

Let the modulator be described by the state recursion
\[
X_{n+1} = A X_n + B \bigl( Q(X_n + U_n) - U_n \bigr),
\]
where \(A\) is an \(L \times L\) matrix corresponding to the \(L\)-fold stack of integrators and feedback coefficients, and \(B\) encodes the structure of the quantizer coupling. The infinite-order limit corresponds formally to letting \(L\to\infty\), although practical modulators are always finite. The purpose of the unlimited-order theory is not to build literal infinite-order modulators but to characterize the generic conditions under which high-order systems remain stable.

The primary observation is that the entropy potential
\[
S(X) = \frac{1}{2} X^\top P X
\]
may be chosen to be uniformly convex independently of the order, provided that the Gramian of the loop filter remains bounded. In classical terms this corresponds to requiring that the feedback coefficients defining the noise transfer function do not grow without bound. In geometric terms it states that the curvature of the entropy manifold remains uniformly positive as the dimension increases.

We therefore introduce the following condition: there exists a constant \(\lambda>0\), independent of the order \(L\), such that the Hessian satisfies
\[
\nabla^2 S(X) \succeq \lambda I
\]
for all states \(X\). Under this uniform convexity assumption we can formulate the unlimited-order stability theorem:

\medskip

\textbf{Theorem A.4.}
\emph{Let \(A\) be the state-transition matrix of an arbitrary-order delta–sigma modulator and let \(B\) represent the quantizer coupling. Suppose the entropy potential \(S\) is uniformly convex, with \(\nabla^2 S \succeq \lambda I\) for some \(\lambda>0\), and that the divergence of the loop filter vector field satisfies}
\[
\mathrm{div}(A X) \le 0
\]
\emph{for all states outside a sufficiently large compact region. Suppose also that the quantization disturbance is uniformly bounded. Then the delta–sigma modulator is globally stable for all bounded inputs, independently of its order.}

\medskip

\textbf{Proof.}
The proof follows directly from the entropy-balance argument developed earlier. Uniform convexity ensures that the entropy potential grows at least linearly with \(\|X\|\), and therefore the entropy gradient provides an unbounded restoring force. The nonpositive divergence of the loop filter vector field implies that the flow does not expand volume in the state space, a condition that prevents the emergence of trajectories that accumulate unbounded energy. Bounded quantization error ensures that the entropy injection never exceeds finite bounds. These three conditions jointly enforce the contractive inequality
\[
S(X_{n+1}) - S(X_n) \le -\alpha S(X_n) + \beta,
\]
with \(\alpha\) independent of modulator order. It follows that the state trajectory must eventually enter a compact sublevel set of the entropy potential and remain within it forever. This proves global stability. \(\square\)

\medskip

This theorem is noteworthy because it reveals that stability is not primarily a question of order but of geometry. The unbounded difficulty of classical stability analysis for high-order modulators arises not from fundamental mathematical barriers but from the representation of the modulator as a linear difference equation perturbed by quantization. The geometric formulation, by contrast, isolates the structural invariants that guarantee stability and shows that these invariants do not deteriorate as the order increases.

\section{A.20 Strong Curvature Bounds and Global Convergence}

Stability is closely related to global convergence of the state trajectories to a bounded attractor. The entropy-based Lyapunov analysis provides sufficient conditions for convergence, but stronger results can be obtained by imposing bounds on the curvature of the entropy landscape. In geometric terms the curvature of the entropy potential dictates the rate at which deviations from equilibrium decay. The stronger the curvature, the faster the convergence.

Let \(K_{\min}\) denote the minimum sectional curvature associated with the metric defined by the Hessian of \(S\). The gradient flow interpretation yields
\[
\frac{d}{dt}\|X(t)\|^2 \le - 2 K_{\min} \|X(t)\|^2 + C,
\]
where \(C\) captures the quantization-induced injection. When \(K_{\min}>0\), Grönwall’s inequality implies the convergence estimate
\[
\|X(t)\|^2 \le e^{-2K_{\min} t} \|X(0)\|^2 + \frac{C}{2K_{\min}} \left(1 - e^{-2K_{\min} t}\right).
\]
The discrete-time recursion inherits this inequality in the limit of small sampling period. This proves that the rate of convergence to the global attractor is determined by the smallest curvature of the entropy manifold.

In classical delta–sigma design this corresponds to selecting loop-filter coefficients so that the noise transfer function has sufficient gain around its high-frequency poles. The geometric formulation shows that these design rules are manifestations of curvature requirements, and that optimizing NTF shape is equivalent to engineering curvature on the entropy manifold.

\section{A.21 Entropy-Driven Control and Nonlinear Quantizers}

Thus far the analysis has focused on uniform quantizers. However, real modulators frequently employ nonuniform quantization, multi-bit quantizers with unequal step sizes, or dynamically adaptively quantizers. The geometric theory accommodates these generalizations naturally by observing that any quantizer can be regarded as a projection onto a discrete submanifold of the state space.

Let \(\Lambda \subset \mathcal{M}\) denote the quantizer image. The projection map \(Q\) defines a retraction of \(\mathcal{M}\) onto \(\Lambda\), and the quantization error is given by \(e(x) = x - Q(x)\). The essential property for stability is that the projection does not increase entropy excessively. This is formalized by the curved-space inequality
\[
S(Q(x)) \le S(x) + \delta(x),
\]
where \(\delta(x)\) depends on the quantizer geometry. For a uniform quantizer \(\delta\) is constant; for a nonuniform quantizer it depends on the spacing of thresholds. The stability condition derived earlier becomes
\[
\Delta S_n = S(X_{n+1}) - S(X_n) \le -\alpha S(X_n) + \sup_x \delta(x).
\]
This inequality shows that stability is preserved so long as the entropy increase produced by the quantizer is uniformly bounded. The nonuniformity of the quantizer does not fundamentally alter the qualitative behavior of the system. Rather, it modifies the shape of the entropy landscape and slightly adjusts the bounds of the global attractor. The geometric theory thus justifies the empirical observation that multi-bit modulators generally exhibit improved stability and lower distortion, since finer quantization reduces the size of the entropy injection.

\section{A.22 Equivalence of Discrete and Continuous Stability Criteria}

The stability analysis can now be unified across the discrete-time and continuous-time formulations. The discrete-time modulator corresponds to a Euler discretization of the entropy-flow PDE, and stability in one domain implies stability in the other provided the sampling period is sufficiently small. The equivalence is formalized in the following theorem.

\medskip

\textbf{Theorem A.5.}
\emph{Let the continuous-time entropy-flow PDE admit a unique global solution that remains bounded for all time, and suppose the entropy potential is uniformly convex. Then the discrete-time delta–sigma modulator obtained by sampling the PDE with sufficiently small sampling period is stable. Conversely, if the discrete-time modulator is stable for all sufficiently small sampling periods, then the continuous-time PDE has a bounded solution.}

\medskip

\textbf{Proof.}
The proof follows from the theory of numerical stability of gradient flows. Uniform convexity of the entropy potential guarantees that the discrete and continuous evolutions remain close in the limit of small sampling period. Boundedness of the continuous solution implies boundedness of its discrete approximation. Conversely, boundedness of the discrete trajectories implies boundedness of the sampled PDE solution through the backward-error interpretation of numerical integration schemes. \(\square\)

\medskip

This theorem shows that delta–sigma stability is not an artifact of discretization but a reflection of the underlying geometry of entropy flows. The sampling operation preserves stability so long as it respects the curvature of the entropy manifold.

\section{A.23 MASH Architectures and Entropy Factorization}

Multi-stage noise-shaping (MASH) modulators provide an important class of delta–sigma structures by cascading lower-order modulators in such a way that their noise contributions cancel. In classical theory this relies on polynomial factorization. In the geometric formulation MASH architectures correspond to a decomposition of the entropy manifold into submanifolds whose curvature properties align to produce cancellation of entropy injections.

Let the state manifold be decomposed into a product \(\mathcal{M} = \mathcal{M}_1 \times \mathcal{M}_2\). The first modulator evolves on \(\mathcal{M}_1\), producing a quantization disturbance whose projection onto \(\mathcal{M}_2\) is input to the second modulator. The cancellation of quantization noise arises when the entropy gradients of the two submanifolds satisfy an orthogonality condition:
\[
\langle \nabla S_1(x), \nabla S_2(y) \rangle = 0.
\]
Under this condition the entropy injection produced by the first stage propagates into directions orthogonal to the entropy gradient of the second stage, and is therefore dissipated harmlessly. The classical polynomial factorization \( (1 - z^{-1})^L = (1 - z^{-1})^{L_1} (1 - z^{-1})^{L_2} \) is thus the algebraic representation of entropy-manifold factorization.

This geometric picture clarifies why MASH structures achieve high-order noise shaping without sacrificing stability: each stage manages its own entropy independently, and the cascading structure aligns their curvature properties so that disturbances from one stage dissipate in the next.

The next section will unify these results in a general curvature-driven characterization of delta–sigma modulators, culminating in the derived-geometric interpretation of quantization as a homotopical projection on a coherence-preserving manifold.

\section{A.24 Curvature–Driven Characterization of Delta–Sigma Architectures}

The preceding developments allow us to approach a unifying description of delta–sigma modulators grounded entirely in the geometry of curvature and entropy routing. Classical architecture diagrams distinguish between feedforward, feedback, cascade, and MASH configurations, but these distinctions obscure the structural invariants that govern stability and noise shaping. From the geometric perspective every delta–sigma modulator is a particular realization of a curvature–constrained flow coupled to a projection map. The curvature dictates the local contraction properties of the entropy manifold, while the projection enforces discrete collapses that maintain coherence by preventing the accumulation of divergence.

Let \(\mathcal{M}\) denote the state manifold endowed with the entropy metric \(g\) derived from the Hessian of the potential \(S\). The loop filter induces a vector field \(v\) whose flows preserve or reduce entropy in the stable region. The quantizer acts as a projection \(Q:\mathcal{M}\to\Lambda\), where \(\Lambda\) is a discrete set representing symbolic output levels. The fundamental behavior of a delta–sigma architecture is encoded in the interaction between three geometric quantities: the sectional curvature \(K(\sigma)\) of the entropy manifold, the divergence \(\mathrm{div}(v)\) of the loop-filter vector field, and the entropy-injection functional induced by quantization.

The curvature \(K(\sigma)\) determines the focusing properties of geodesics on \(\mathcal{M}\). When curvature is strictly positive, geodesics converge and the gradient flow tends toward an attractor. When curvature is zero, geodesics neither converge nor diverge, and stability depends on additional constraints from the vector field. When curvature is negative, geodesics diverge and the system risks instability unless counteracted by quantization-induced projection. The curvature signature therefore characterizes the inherent stability of the loop filter independent of the quantizer. Classical design rules that place zeros of the noise transfer function at the origin can be interpreted as engineering positive curvature in the low-frequency subspace of the manifold.

The divergence \(\mathrm{div}(v)\) quantifies how the loop filter compresses or expands volume in the state space. In delta–sigma modulators the divergence is often negative outside small neighborhoods of the origin, reflecting the dissipative nature of the integrator structure. This corresponds to classical pole-placement rules that ensure the effective poles lie within the unit circle. The geometric interpretation shows that these rules enforce a global volume contraction that guarantees the existence of an entropy well.

Finally, the quantizer projection introduces discrete curvature corrections that counteract any residual divergence. These corrections are represented by impulses in the entropy-flow PDE and ensure that the system does not drift along neutral directions of the manifold. Delta–sigma architectures differ in how they combine these curvature contributions, but they all operate according to the same principles. Thus the architecture of a modulator is fully described by the triplet \((\mathcal{M},v,Q)\), and all classical configurations are merely coordinate realizations of this geometric object.

\section{A.25 Homotopical Quantization as Derived Projection}

Quantization becomes nontrivial at higher orders because the projection map \(Q\) imposes constraints that do not preserve smoothness. These constraints introduce discontinuities in the evolution of the state, forcing the geometric structure of the manifold to adapt instantaneously. In classical approaches these nonsmooth effects are treated as artifacts of discretization. The derived-geometric interpretation shows that they are essential structural operations analogous to truncations in homotopy theory.

Let the quantizer image \(\Lambda\) be regarded as a discrete derived subspace of the manifold \(\mathcal{M}\). The projection \(Q\) is then a derived morphism, meaning that it preserves not only the pointwise structure but also the higher-order homotopical relations encoded in the derived tangent complex. The action of the quantizer on the state amounts to a homotopy that replaces the full tangent complex at a given point by a truncated version compatible with membership in \(\Lambda\). This truncation removes unstable components of the tangent directions that would otherwise lead to unbounded excursions. 

The derived projection property can be expressed formally by the condition
\[
\mathrm{fib}(X \to \Lambda) \simeq T_X\mathcal{M}_{\leq 0},
\]
where the right-hand side denotes the connective part of the derived tangent complex. The quantizer enforces that the state must evolve along directions whose curvature and entropy characteristics satisfy certain coherence constraints. The nonlinearity of quantization thus corresponds to a higher-order correction to the geometry of the state manifold, and the stability of the modulator arises from the compatibility of these truncations with the curvature structure of \(\mathcal{M}\).

This perspective shows that quantization is not merely the introduction of error but the enforcement of compatibility conditions between the state and the manifold within which the state evolves. As such, quantization is a homotopical correction that ensures the persistence of coherence under perturbation. Stability emerges because the projection eliminates directions in which entropy would increase without bound.

\section{A.26 Free–Energy Equivalence and Generalized Lyapunov Functionals}

The interpretation of delta–sigma modulation as a gradient flow perturbed by impulsive forcing invites a connection to generalized free–energy principles. In this context the entropy potential \(S\) plays a role analogous to a free–energy functional, and stability corresponds to the existence of a descent direction for \(S\) along the modulator trajectories. Classical Lyapunov functions represent special instances of such functionals. The geometric viewpoint reveals that free–energy functionals arise naturally from the curvature of the entropy manifold, and that their descent structure governs stability.

Let \(F(X) = S(X) + R(X)\) be a generalized free–energy functional in which \(R(X)\) captures additional structure introduced by the loop filter. The modulator update can then be written as
\[
X_{n+1} = \mathrm{prox}_{\Delta t F}(X_n) + B e_n,
\]
where the proximal map is defined implicitly by the variational condition
\[
X_{n+1} = \arg\min_Y \left( F(Y) + \frac{1}{2\Delta t}\|Y - X_n\|^2\right).
\]
In the absence of quantization the proximal map generates a descent in the free–energy functional. Quantization introduces a bounded perturbation, which can be interpreted as a bounded additive disturbance in the variational problem. Stability follows when the dissipation from the free–energy descent dominates the entropy injection from quantization.

Furthermore, by choosing \(R\) appropriately, one can express classical stability criteria such as the root conditions on the loop filter in terms of the curvature of \(F\). The general criterion for stability becomes: the free–energy functional must be convex and must decrease under the proximal flow. Quantization must not overwhelm the dissipative structure. This generalizes all classical Lyapunov stability proofs and reveals the underlying variational principle that governs delta–sigma modulators.

\section{A.27 Geometric Noise–Shaping Limits and Optimal Curvature}

The noise–shaping performance of a delta–sigma modulator is intimately tied to the curvature of the entropy manifold. High-order modulators exhibit steep suppression of quantization noise at low frequencies, a property that is reflected geometrically as high curvature in regions where the gradient flow is most active. The optimality of a modulator in terms of noise-shaping therefore corresponds to optimizing the curvature profile of the entropy potential.

Let the noise spectral density be denoted \(N(\omega)\). Classical theory shows that in an \(L\)-th order modulator the in-band noise scales as \(\omega^{2L}\). In the geometric formulation this scaling law corresponds to the statement that the entropy manifold exhibits curvature of order \(L\) in the low-frequency subspace. More formally, if \(\sigma_\omega\) is the two-plane corresponding to a frequency mode \(\omega\), then
\[
K(\sigma_\omega) \sim \omega^{2L}
\]
for small \(\omega\). The curvature thus determines the growth rate of the noise transfer function near the origin.

The noise–shaping limit arises because curvature cannot increase without bound. Excessive curvature leads to numerical stiffness in the gradient flow, making the system sensitive to quantization disturbances and prone to overload. The optimal noise–shaping tradeoff corresponds to selecting the highest curvature consistent with the convexity of the entropy potential and the boundedness of quantization injection. In classical terms this is the well-known tradeoff between high-order noise shaping and reduced stability margin. The geometric theory provides a unified explanation: stability fails when curvature exceeds the level at which the truncated derived tangent complex can enforce coherence.

\section{A.28 Toward a Unified Curvature Criterion for Stability and Performance}

The results accumulated throughout this appendix suggest the existence of a unified criterion that governs both stability and noise-shaping performance in delta–sigma modulation. This criterion takes the form of a bound on the curvature of the entropy manifold. Stability requires that curvature be positive and bounded away from zero. Noise-shaping performance demands that curvature increase rapidly near low-frequency modes. The optimal modulator is therefore one in which the entropy manifold possesses a carefully tuned curvature profile.

We express this criterion formally as follows: let the sectional curvature \(K(\sigma)\) of the entropy manifold satisfy
\[
0 < K_{\min} \le K(\sigma) \le K_{\max} < \infty,
\]
where \(K_{\min}\) determines stability and \(K_{\max}\) determines performance. Excessively low curvature leads to insufficient dissipation of quantization errors, while excessively high curvature leads to instability. The ratio \(K_{\max}/K_{\min}\) thus becomes a fundamental measure of the modulator’s performance–stability tradeoff. Classical NTF design rules correspond to selecting feedback coefficients to engineer this ratio.

The derived-geometric interpretation closes the loop by showing that quantization enforces a truncation of the tangent complex that ensures that curvature remains bounded above. This truncation is the homotopical constraint that stabilizes the system in the presence of curvature-driven noise shaping. The next sections will extend these results into a complete variational characterization of delta–sigma modulation, establishing the formal equivalence between quantization, curvature, free–energy descent, and coherence–preserving homotopy.

\section{A.29 Variational Characterization of Delta–Sigma Modulation}

The geometric framework developed so far points toward a fully variational interpretation of delta–sigma modulation. Classical treatments view the modulator as a feedback system whose purpose is to linearize quantization and shape noise out of the signal band. The variational interpretation reveals that this behavior emerges from a deeper principle: the modulator performs implicit minimization of a free–energy functional defined on the entropy manifold. Quantization, feedback, and noise shaping become natural consequences of the structure of this functional and the constraints imposed by discrete projection.

Let the free–energy functional be written as
\[
\mathcal{F}(X) = S(X) + R(X) - \langle X, U \rangle,
\]
where \(S(X)\) is the entropy potential, \(R(X)\) encodes the loop-filter geometry, and the final term incorporates the signal input. The negative sign arises from the fact that the modulator attempts to match its internal state to the input by minimizing mismatch energy. The evolution of the modulator can be interpreted as a discrete approximation to the gradient flow of \(\mathcal{F}\),
\[
X_{n+1} = X_n - \Delta t \, \nabla \mathcal{F}(X_n) + B e_n,
\]
where the term \(B e_n\) represents quantization-induced impulses. In the absence of quantization the system would converge to the minimizer of \(\mathcal{F}\). The presence of quantization introduces bounded disturbances, and the structure of the manifold ensures that the system remains in the vicinity of the minimizer while exporting quantization noise into directions orthogonal to the input.

This variational formulation clarifies why delta–sigma modulators exhibit robust linearity despite the nonlinearity of quantization. The minimization of the generalized free energy forces the state trajectory to encode the input signal in a latent representation where the quantization disturbance appears orthogonal to the primary descent direction. Classical NTF behavior emerges because the descent in \(\mathcal{F}\) is suppressed in certain directions and promoted in others. The zero structure imposed by the loop filter corresponds to flat or saddle-like regions of the free–energy landscape. Thus the entire modulator architecture is encoded in the geometry of \(\mathcal{F}\), and stability becomes the consequence of its convexity properties.

\section{A.30 Homotopy, Coherence, and Derived Lyapunov Structure}

The Lyapunov function used in classical stability proofs is a special case of a more general derived structure. In the context of derived geometry the Lyapunov function corresponds not merely to a scalar potential but to a homotopy class of higher-order structures that preserve coherence under perturbations. The quantizer collapses the state onto a discrete subset of the manifold, and the Lyapunov function must remain compatible with the truncated tangent directions at all points in \(\Lambda\).

Let the derived tangent complex at a point \(X\in\mathcal{M}\) be written as
\[
T_X^\bullet \mathcal{M} = \left( T_X \mathcal{M} \to L_X \right),
\]
where \(L_X\) encodes the relations and obstructions associated with the derived structure. Stability requires that the differential of the Lyapunov function annihilate all components of \(T_X^\bullet\mathcal{M}\) that correspond to unstable directions. Formally this requires
\[
dS(X)\bigl(v_i\bigr) \le 0
\]
for every vector \(v_i\) in the nonpositive-degree truncation of the derived tangent complex. This ensures that the gradient of the Lyapunov potential always points inward along directions capable of producing instability.

Quantization modifies the tangent complex by projecting onto a derived subspace. The Lyapunov function must be compatible with this projection, meaning that the difference
\[
S(Q(X)) - S(X)
\]
must remain bounded above. This bound ensures that quantization never injects excessive entropy. The derived Lyapunov condition can thus be written as
\[
S(X_{n+1}) - S(X_n) \le -\alpha S(X_n) + \beta,
\]
which incorporates both the variational descent and the homotopical correction introduced by quantization. This condition encapsulates the entire stability theory of delta–sigma modulators and represents its most general homotopical form.

\section{A.31 Geometric Limits: Overload, Surge, and Collapse}

Although delta–sigma modulators are designed to be stable under normal operation, they can undergo instability when the input amplitude exceeds certain bounds or when the internal state is perturbed into certain high-energy regions. This behavior is traditionally described as overload. The geometric framework explains overload as the crossing of certain curvature thresholds beyond which the entropy manifold can no longer enforce contraction of the state trajectory.

Let \(X^\ast\) denote the equilibrium state corresponding to a given input. Overload occurs when the modulator state leaves the stable curvature region around this equilibrium, entering a region where the entropy potential becomes locally concave. In such regions the gradient flow no longer points toward the equilibrium, and quantization disturbances may increase entropy rather than decrease it. This phenomenon corresponds to negative curvature in the entropy manifold.

Surge refers to transient excursions in which the state moves rapidly through regions of low curvature before returning to the stable region. Such surges are common in practical modulators and do not necessarily indicate instability. They arise because the curvature profile of the entropy manifold is not uniform, and certain directions may temporarily exhibit weak focusing properties. The derived tangent complex ensures that the quantizer forces the state back into the stable region once the surge subsides.

Collapse refers to a more severe phenomenon in which the state enters a region of the manifold where curvature is insufficient to counteract divergence, and the quantizer cannot supply enough corrective impulse to re-establish coherence. In this scenario the state trajectory escapes to infinity or oscillates with unbounded amplitude. Classical designers interpret this as catastrophic instability. The geometric theory interprets it as a topological obstruction: the quantizer projection fails to induce the necessary truncation of the tangent complex because the curvature conditions have been violated.

These concepts reveal that overload, surge, and collapse are not independent phenomena but different manifestations of curvature-driven transitions in the entropy manifold. Their analysis provides deeper insight into robustness and informs new design methodologies based on curvature shaping.

\section{A.32 RSVP Integration: Scalar–Vector–Entropy Decomposition}

The RSVP theoretical framework introduces an alternative decomposition of fields into scalar, vector, and entropy components. This decomposition applies naturally to the state-space description of delta–sigma modulators. The scalar component corresponds to the entropy potential \(S(X)\), the vector component corresponds to the loop filter vector field \(v(X)\), and the entropy component corresponds to the quantization-induced impulses. The scalar–vector–entropy decomposition thus provides a complete representation of the forces acting on the modulator state.

Let the decomposition be written as
\[
\dot{X} = -\nabla S(X) + v(X) + \eta(t),
\]
where \(\eta(t)\) represents quantization impulses. The stability of the modulator depends on the balance between the scalar gradient, which enforces entropy descent, and the vector flow, which redistributes energy across the manifold. The quantization impulses serve as entropy injections that allow the system to maintain coherence by correcting the discretization-induced deviations.

The scalar–vector–entropy decomposition reveals that delta–sigma modulation is a specific instance of a broader class of systems governed by entropy shaping and feedback. This aligns delta–sigma theory with the general principles of RSVP, which describe how systems maintain coherence by redirecting entropy along constrained geometric flows. The modulator is therefore a microcosm of the RSVP dynamics, and the stability theory developed here reflects the same balance of forces seen in larger physical and cognitive systems.

\section{A.33 Derived Mapping Stacks and Measurement Coherence}

The measurement performed by a delta–sigma modulator maps a continuous-time signal into a discrete bitstream. This mapping is represented geometrically as a morphism between a smooth manifold and a discrete set. In classical signal processing the measurement is viewed as a sampling and quantization operation. The derived-geometric approach shows that the mapping is more properly viewed as a morphism in the category of derived stacks.

Let the input signal space be a smooth derived stack \(\mathcal{X}\), and let the output bitstream space \(\mathcal{Y}\) be a discrete derived stack. The delta–sigma modulator implements a derived morphism
\[
f : \mathcal{X} \longrightarrow \mathcal{Y},
\]
subject to coherence constraints imposed by the entropy-flow geometry. The derived tangent complex of this morphism encodes the quantization constraints. Stability corresponds to the existence of a lift of the morphism to a homotopically coherent diagram in which truncations preserve the compatibility of the tangent complexes. In this interpretation delta–sigma modulation is a special case of active measurement on derived stacks.

This framework unifies the physical, informational, and geometric perspectives on delta–sigma modulation. The bitstream is not simply a representation of the input; it is the result of a homotopically coherent mapping in which entropy is routed along stable geometric directions. The stability proofs in preceding sections reveal the structural conditions under which this mapping preserves coherence.

\section{A.34 Limits of Derived Coherence and Quantization Obstructions}

The derived–geometric interpretation of delta–sigma modulation identifies stability as the persistence of coherence in the presence of quantization-induced discontinuities. Coherence is maintained so long as the derived tangent complex of the measurement mapping remains compatible with the curvature conditions of the entropy manifold. There exist, however, intrinsic limits to this compatibility that arise from obstructions in the derived structure. These obstructions define the ultimate bounds of delta–sigma performance, independent of classical linear analysis.

Let the quantizer projection \(Q\) be viewed as a derived morphism from the state manifold \(\mathcal{M}\) to the discrete output space \(\Lambda\). The derived tangent complex at a point \(X\in\mathcal{M}\) maps to a truncated tangent complex at the image \(Q(X)\). This mapping preserves coherence only if the derived differential
\[
dQ : T_X^\bullet\mathcal{M} \longrightarrow T^\bullet_{Q(X)}\Lambda
\]
is homotopically surjective onto the connective part of the target complex. When this surjectivity fails, quantization injects a distortion that cannot be corrected by entropy descent or vector-field dissipation. The obstruction to surjectivity manifests as an uncorrectable mode in the state trajectory, corresponding to a direction along which entropy increases under the flow. Classical descriptions of this phenomenon refer to it as noise folding or high-order instability, but the derived framework reveals that it is fundamentally a failure of homotopical coherence.

The quantization obstruction is therefore encoded in a derived cohomology class
\[
\mathrm{Obs}(Q) \in H^1\left(T^\bullet\mathcal{M}, T^\bullet\Lambda \right),
\]
which vanishes exactly when coherence is preserved. In practice, this cohomology class measures the incompatibility between the curvature of the entropy manifold and the topology of the quantizer image. Stability requires that the obstruction class vanish uniformly across the entire stable region. When the input amplitude grows or the loop-filter coefficients push the curvature beyond admissible bounds, the obstruction class becomes nontrivial and the modulator loses coherence. This provides a rigorous geometric explanation for the empirical phenomenon of overload-induced instability.

\section{A.35 Global Invariants and the Entropy–Curvature Product}

The stability of a delta–sigma modulator is governed not only by local curvature conditions but also by global invariants of the entropy manifold. These invariants determine how the state trajectory explores the manifold and whether the combination of curvature and quantization supports long-term bounded evolution. The principal invariant is the entropy–curvature product, which captures the interplay between curvature-induced contraction and quantization-induced entropy injection.

Let the sectional curvature be denoted \(K(\sigma)\) and the quantization injection at state \(X\) be denoted \(\Delta S_Q(X)\). The entropy–curvature product is defined formally as
\[
\Gamma = \sup_{X \in \mathcal{M}} \left( \frac{\Delta S_Q(X)}{K(\sigma_X)} \right),
\]
where \(\sigma_X\) denotes the two-plane in the tangent space associated with the direction of quantization error. Stability requires that \(\Gamma\) be finite. If \(\Gamma\) diverges, quantization injects entropy faster than curvature can dissipate it, leading to runaway trajectories even when local convexity is present.

The invariant \(\Gamma\) plays a role analogous to the spectral radius of the classical noise transfer function. Classical NTF design ensures that the magnitude of the NTF remains within certain bounds to avoid excessive amplification of quantization noise. The geometric invariant \(\Gamma\) expresses the same principle in intrinsic terms: the curvature of the manifold must be sufficiently strong to counterbalance the worst-case entropy injection. The classical condition that poles of the loop filter lie within the unit circle corresponds to ensuring that the curvature does not flatten in any direction of the manifold.

Global invariants also govern the maximum achievable noise-shaping order. As curvature increases to achieve higher-order suppression of low-frequency noise, the entropic cost per quantization step increases as well. At some point curvature scaling outpaces the ability of entropy descent to stabilize the system. This yields a geometric explanation for the classical empirical observation that modulators beyond eighth or ninth order become increasingly difficult to stabilize. The difficulty arises not from algebraic complexity but from topological constraints on the entropy–curvature product.

\section{A.36 Complete Characterization of Curvature–Driven NTF Design}

Noise shaping in delta–sigma modulation is traditionally implemented through the design of the noise transfer function. The geometric theory reveals that the NTF can be interpreted as the signature of the curvature of the entropy manifold along the directions orthogonal to the signal. This relationship makes explicit the correspondence between pole–zero placement and curvature engineering.

Let the NTF be represented in the classical form
\[
\mathrm{NTF}(z) = \prod_{i=1}^L (1 - z^{-1}) p_i(z),
\]
where the \(p_i(z)\) correspond to the classical loop-filter factors. In the geometric framework the roots at \(z = 1\) introduce directions of extremely high curvature near low-frequency modes. The multiplicity of the root determines the order of the curvature in the low-frequency subspace. The additional factors \(p_i(z)\) shape the curvature in higher-frequency regions.

The curvature tensor of the entropy manifold may be expanded in a local coordinate chart as
\[
R_{ijkl}(X) = \sum_{m=1}^L c_m(X) \, \omega_{ijkl}^{(m)},
\]
where the coefficients \(c_m(X)\) correspond to curvature contributions of order \(m\), and the tensors \(\omega_{ijkl}^{(m)}\) represent the geometric directions shaped by the modulator structure. The mapping between NTF factors and curvature components is then given by
\[
c_m(X) \propto \frac{d^m}{dz^m} \mathrm{NTF}(z)\bigg|_{z=1}.
\]
This identifies the curvature of order \(m\) with the \(m\)-th derivative of the NTF at the origin of its frequency domain. Thus NTF design becomes a matter of tailoring the curvature spectrum of the entropy manifold. Stability corresponds to ensuring that the curvature does not exceed the bounds imposed by the derived tangent truncation.

This perspective unifies all classical NTF constraints, including Lee’s criterion and the crest-factor bounds, under a single geometric principle: the curvature spectrum must remain compatible with quantization-induced entropy injection. Excessive curvature in any direction overwhelms the truncation mechanism, leading to instability. Optimal design balances curvature across all directions, yielding maximum noise shaping within the limits of derived coherence.

\section{A.37 Noise Shaping as a Derived PDE System}

Noise shaping can also be interpreted as the solution of a derived partial differential equation whose structure reflects the interplay between entropy, curvature, and quantization. The continuous-time analogue of delta–sigma modulation obeys an entropy-flow PDE of the form
\[
\frac{\partial X}{\partial t} = -\nabla S(X) + v(X) + \eta(t),
\]
where \(\eta(t)\) encodes quantization impulses. The discrete-time modulator corresponds to a sampling and projection of this PDE onto the quantizer image. The noise-shaping behavior emerges as a property of the PDE solution’s spectral decomposition.

Let the manifold be equipped with a Laplace–Beltrami operator \(\Delta_g\) corresponding to the entropy metric. The linearization of the PDE around an equilibrium yields
\[
\frac{\partial}{\partial t} \delta X = - H_S \delta X + J_v \delta X + \eta(t),
\]
where \(H_S\) is the Hessian of the entropy potential and \(J_v\) is the Jacobian of the vector field. The spectral decomposition of the operator \(-H_S + J_v\) yields a natural frequency basis for the manifold. The corresponding eigenvalues control the rate at which disturbances decay or grow.

Noise shaping occurs because the eigenvalues associated with low-frequency modes are made exceedingly large in magnitude, causing quantization disturbances to decay rapidly in these directions while being allowed to propagate in orthogonal directions. In effect, noise shaping is equivalent to a curvature-weighted diffusion process. The NTF represents the discrete-time analogue of the Green’s function of this derived PDE system.

By expressing noise shaping as a PDE problem, the derived-geometric theory provides a unified account of classical frequency-domain behavior. The magnitude of the NTF at any frequency reflects the amount of curvature-induced dissipation in the corresponding eigenmode of the Laplace–Beltrami operator. This PDE perspective also generalizes naturally to multi-dimensional and MIMO delta–sigma modulators, revealing that the same principles govern noise shaping in spatially coupled systems.

\section{A.38 Derived Curvature Bounds and the Maximum Stable Order}

The preceding analysis implies that the maximum stable order of a delta–sigma modulator can be expressed as a bound on derived curvature. Let the curvature spectrum of the entropy manifold be denoted by \(\{K_m\}\), where \(K_m\) represents the curvature contribution of order \(m\). Stability requires that the derived curvature satisfy
\[
K_m \le K_{\mathrm{max}}
\]
for all orders \(m\), where \(K_{\mathrm{max}}\) is the largest curvature compatible with the quantizer truncation. This bound arises because the quantizer can only eliminate unstable components of the tangent complex up to a finite degree of nonlinearity. Beyond this degree the curvature overwhelms the projection, leading to loss of coherence.

The largest order \(L_{\mathrm{max}}\) for which the modulator remains stable satisfies
\[
L_{\mathrm{max}} = \max \left\{ m \, \middle| \, K_m \le K_{\mathrm{max}} \right\}.
\]
Classical results indicate that modulators of order greater than about eight exhibit practical instability. The derived–geometric framework explains this empirical limit by showing that the curvature spectrum of typical modulator architectures grows superlinearly with order, eventually violating the coherence bound.

This curvature-based order limit applies universally, regardless of whether the modulator is implemented in analog, digital, switched-capacitor, or MASH form. The derived curvature bound therefore represents the first known architecture-independent explanation for the practical upper limit of delta–sigma order. The next sections shall expand these curvature bounds into a complete geometric–variational characterization of performance and stability.

\section{A.39 Global Spectral Geometry and the Delta–Sigma Attractor}

The spectral properties of the entropy manifold determine not only the local stability and noise–shaping behavior of the modulator but also the nature of its global attractor. The attractor represents the set of states toward which the system converges in the presence of bounded quantization impulses. It serves as a geometric encoding of the modulator’s long–term behavior and influences everything from distortion characteristics to overload recovery. Classical analysis describes the attractor as a bounded invariant set of the difference equation. The geometric framework reveals a deeper structure: the attractor is a compact subset of the entropy manifold endowed with spectral properties inherited from the Laplace–Beltrami operator of the entropy metric.

Let the Laplace–Beltrami operator be denoted \(\Delta_g\). Its eigenfunctions \(\phi_k\) and eigenvalues \(\lambda_k\) govern the diffusion of entropy on the manifold and thus determine how perturbations decompose across different geometric directions. The attractor \(\mathcal{A}\) is characterized by the property that it is the minimal compact set invariant under the discrete flow
\[
X_{n+1} = f(X_n) + B e_n
\]
and that its tangent directions correspond exactly to eigenfunctions associated with nonpositive eigenvalues of the linearized operator. Positive eigenvalues correspond to unstable directions, but these are removed by quantizer truncation, leaving only the stable portion of the spectrum.

Thus the attractor can be characterized as the image of the stable spectral projection
\[
\Pi_{\mathrm{stable}} : T_X\mathcal{M} \longrightarrow T_X\mathcal{A},
\]
where \(T_X\mathcal{A}\) denotes the tangent space of the attractor at the point \(X\). This reveals that the attractor inherits geometric structure from the manifold: its curvature, topology, and dimensionality reflect the stable spectrum of \(\Delta_g\). The attractor is therefore not an arbitrary shape but a geometric object constrained by the interplay between curvature, divergence, and quantization.

In particular, the dimension of the attractor can be approximated by the count of eigenvalues \(\lambda_k\) below a threshold determined by the entropy injection. This provides a geometric explanation for the classical observation that the attractor of a delta–sigma modulator often appears low–dimensional, even in high–order systems. The attractor dimension is constrained by the number of geometric directions in which curvature remains sufficiently strong to dissipate quantization disturbances. The derived–geometric analysis thus yields a spectral characterization of the attractor and ties its structure directly to stability and performance.

\section{A.40 Homotopical Quantization and Symbolic Collapse}

The discrete output of a delta–sigma modulator can be interpreted as a symbolic representation of the input. Classical treatments describe quantization as the conversion of continuous states into discrete symbols. The homotopical interpretation reveals a richer structure in which quantization corresponds to a collapse of the derived tangent complex. This collapse is neither arbitrary nor purely combinatorial. It is governed by coherence constraints that ensure the compatibility of symbolic states with the underlying geometry of the manifold.

Let the derived tangent complex at a point \(X \in \mathcal{M}\) be written as
\[
T_X^\bullet = \left( V_0 \to V_1 \right),
\]
where \(V_0\) is the classical tangent space and \(V_1\) encodes obstructions and higher–order relations. The quantizer projection \(Q\) induces a truncation
\[
T_{Q(X)}^\bullet = \tau_{\le 0} T_X^\bullet,
\]
which removes higher–order relations and retains only the connective portion of the complex. This truncation is called symbolic collapse. It represents the fundamental act of discretization: the system collapses the full geometric structure of the signal onto a finite symbolic alphabet.

Coherence requires that the truncated structure remain compatible with the curvature of the entropy manifold. If curvature is too large in a given direction, the truncation discards essential geometric information and the system loses stability. If curvature is too small, the state drifts through a region where the quantizer does not provide sufficient correction and coherence is again lost. Stability therefore requires that curvature and truncation be matched in a delicate balance, a principle absent from classical descriptions but central to the derived–geometric theory.

Symbolic collapse also plays a role in shaping the error spectrum of the modulator. The error sequence \(\varepsilon_n = X_n - Q(X_n)\) is governed by the mismatch between the full tangent complex and its truncation. The derived structure shows that this mismatch is orthogonal to the stable directions of the entropy manifold, which explains the classical phenomenon that quantization noise is largely uncorrelated with the input signal. Thus symbolic collapse is the geometric origin of white or shaped quantization noise.

\section{A.41 Thermodynamic Limits: Entropy Production and Irreversibility}

Delta–sigma modulators have long been known to exhibit thermodynamic behavior. Their structure resembles feedback–driven dissipative systems that consume entropy locally and expel it globally. The derived–geometric theory provides a foundation for formalizing these thermodynamic properties.

The entropy potential \(S(X)\) governs the dissipation of quantization disturbances. Each quantizer event introduces a finite entropy increment \(\Delta S_Q\), and the gradient flow dissipates this increment over time. The quantization cycle therefore acts as a thermodynamic engine. Let the average entropy injection per cycle be
\[
\langle \Delta S_{\mathrm{inj}} \rangle,
\]
and the average entropy dissipated by the gradient flow be
\[
\langle \Delta S_{\mathrm{diss}} \rangle.
\]
Stability requires the inequality
\[
\langle \Delta S_{\mathrm{diss}} \rangle \ge \langle \Delta S_{\mathrm{inj}} \rangle.
\]
When equality holds, the modulator operates at its thermodynamic limit, achieving maximum noise–shaping performance. When the injection exceeds the dissipation, entropy accumulates in the state trajectory, driving it toward overload and collapse.

Irreversibility arises because quantization discards information that cannot be recovered. The projection \(Q(X)\) is not injective, and its fibers correspond to equivalence classes of states whose differences cannot be reconstructed from the output. This loss of information corresponds to entropy production in thermodynamic systems. The rate of entropy production is given by
\[
\sigma = \frac{\langle \Delta S_{\mathrm{inj}} \rangle - \langle \Delta S_{\mathrm{diss}} \rangle}{\Delta t}.
\]
Stability requires \(\sigma \le 0\). The boundary case \(\sigma = 0\) corresponds to a perfectly balanced modulator, operating as an entropy–neutral engine.

The thermodynamic characterization reveals that delta–sigma modulation is an instance of a broader class of irreversible feedback systems. These systems maintain coherence by dissipating entropy while performing symbolic projection. The geometric theory thus aligns delta–sigma modulators with thermodynamic engines, active inference systems, and biological regulatory loops.

\section{A.42 Entropy–Geometric Equilibrium and the Structure of Steady–State Dynamics}

When the entropy injection and dissipation reach equilibrium, the delta–sigma modulator enters a steady–state regime. This regime corresponds to the global attractor discussed in Section A.39 and represents the state of long–term operation under bounded inputs. The structure of this equilibrium reflects the geometry of the entropy manifold, the shape of the free–energy functional, and the truncation induced by quantization.

Let the steady–state trajectory be denoted \(X_n^\ast\). The equilibrium condition requires
\[
\mathbb{E}\bigl[ S(X_{n+1}^\ast) - S(X_n^\ast) \bigr] = 0,
\]
which implies
\[
\mathbb{E}\bigl[ -\nabla S(X_n^\ast) \cdot v(X_n^\ast) \bigr] + \mathbb{E}\bigl[ \Delta S_Q(X_n^\ast) \bigr] = 0.
\]
This equation expresses the balance between entropy descent along the vector field and entropy injection from quantization. The first term represents geometric dissipation, while the second term represents symbolic injection. Their equality defines the equilibrium structure.

The form of this equilibrium reveals the modulator’s deep connection to active inference. In active inference systems the equilibrium arises from a balance between prediction–error minimization and stochastic fluctuations. In delta–sigma modulators prediction corresponds to the loop filter, error corresponds to quantization noise, and fluctuations correspond to the symbolic collapse. Thus delta–sigma modulation can be viewed as a primitive active–inference system operating on a geometric manifold. Its steady–state dynamics encode the resolution of competing forces: geometric dissipation and symbolic discreteness.

The next sections will deepen this connection by examining the variational structure underlying the equilibrium and providing curvature–dependent estimates for the attractor geometry.

\section{A.43 Variational Equilibria and Curvature–Dependent Bounds}

The equilibrium of a delta–sigma modulator arises from a delicate interaction between the variational structure of the free–energy functional, the curvature of the entropy manifold, and the discrete collapses induced by quantization. Classical descriptions of steady–state behavior as the formation of a “limit cycle” fail to capture the rich geometry underlying this equilibrium. The derived–geometric formulation reveals that the steady state is a variational equilibrium, defined not by a fixed point of the dynamical system but by a balance of geometric and symbolic forces.

Let the full dynamics be expressed as
\[
X_{n+1} = X_n - \Delta t \, \nabla \mathcal{F}(X_n) + B e_n,
\]
where \(\mathcal{F}\) is the generalized free–energy functional defined earlier. The equilibrium condition requires that the expected descent of \(\mathcal{F}\) vanish:
\[
\mathbb{E}\bigl[ \mathcal{F}(X_{n+1}) - \mathcal{F}(X_n) \bigr] = 0.
\]
This is not identical to requiring that the gradient vanish. Instead, equilibrium emerges when the variational descent induced by \(-\nabla \mathcal{F}\) is exactly balanced by the entropy injection from the quantizer. This balance defines a variational fixed point of the modulator, even though the state itself continues to move.

The presence of curvature modifies the geometry of this variational equilibrium. In particular, the second derivative of \(\mathcal{F}\), given by the sum of the Hessians of the entropy potential and the loop–filter contribution, determines the stability of the equilibrium. Let the Hessian be written as
\[
H_{\mathcal{F}}(X) = H_S(X) + H_R(X).
\]
The curvature–dependent bound on perturbations around equilibrium is given by
\[
\| \delta X_{n+1} \| \le \| (I - \Delta t \, H_{\mathcal{F}}(X^\ast)) \delta X_n \| + C,
\]
where \(C\) is the bound on quantization injection. The spectral radius of \(I - \Delta t \, H_{\mathcal{F}}(X^\ast)\) must be less than unity in the stable region. This condition yields the curvature bound
\[
\lambda_{\min}(H_{\mathcal{F}}) > 0,
\]
where \(\lambda_{\min}\) denotes the smallest eigenvalue of the Hessian. When curvature becomes too small in any direction, the variational equilibrium destabilizes. When curvature becomes too large, quantization cannot preserve compatibility and coherence is lost through symbolic collapse that destroys the tangent structure needed for stability. Thus, the variational equilibrium exists only within a bounded curvature region, determined by the geometry of \(\mathcal{F}\) and the truncation behavior of the quantizer.

\section{A.44 Generalized Attractor Measures and Coherent Output Geometry}

The steady–state behavior of a delta–sigma modulator is encoded not only in the trajectory of the internal state but also in the distribution of output symbols. Classical descriptions refer to this distribution as the limit–cycle pattern or the steady–state bitstream. The derived–geometric interpretation reveals that the output bitstream is the observable imprint of the attractor measure on the quantizer image. This measure is induced by the invariant measure of the entropy–flow system under symbolic collapse.

Let \(\mu\) denote the invariant measure supported on the attractor \(\mathcal{A}\). The output measure \(\nu\) on the quantizer image \(\Lambda\) is defined by the pushforward
\[
\nu = Q_\ast \mu.
\]
The geometry of the attractor determines the structure of \(\nu\). Regions of the attractor with high curvature concentrate the measure, producing long runs of repeated symbols in the output. Regions of low curvature disperse the measure, producing frequent symbol alternation. The bitstream is therefore a symbolic representation of the geometric properties of the attractor.

The invariance of \(\mu\) under the derived flow requires that it satisfy the Frobenius–Perron equation associated with the nonlinear transformation \(X \mapsto X - \Delta t \nabla \mathcal{F}(X) + B e\). The existence of a unique attractor measure follows from the contraction properties derived earlier. The regularity of this measure depends on curvature: when curvature is strictly positive and bounded, the attractor measure is absolutely continuous with respect to the Riemannian volume measure on \(\mathcal{M}\). When curvature approaches zero in certain directions, singular measures arise, corresponding to pathological resonance conditions and limit cycles.

From the signal processing perspective the attractor measure determines critical performance metrics such as harmonic distortion, idle tones, and stochastic dithering properties. The derived–geometric framework thus offers a natural explanation for these classical phenomena. Idle tones correspond to singularities in the attractor measure, arising from curvature degeneracies. Dither eliminates such singularities by smoothing the attractor measure through stochastic dispersion. This smoothing is equivalent to modifying the derived tangent complex by injecting additional degrees of freedom, which remove coherence–destroying degeneracies.

Thus, the output geometry of the modulator is a direct manifestation of the geometry of the attractor, and the attractor is itself governed by curvature–dependent variational equilibria.

\section{A.45 Formal Equivalence Between Delta–Sigma Dynamics and Active–Inference Flows}

The connection between delta–sigma modulation and active inference becomes precise when viewed through the lens of variational geometry. Active inference systems minimize a free–energy functional that balances prediction accuracy and uncertainty. Delta–sigma modulators minimize a generalized free–energy functional that balances signal tracking and quantization–induced entropy. The mathematical structure underlying both processes is identical.

In active inference, the state evolves according to
\[
\dot{X} = -\nabla F(X,U) + \xi(t),
\]
where \(F\) is a free–energy functional and \(\xi(t)\) represents stochastic fluctuations. In delta–sigma modulation the corresponding equation is
\[
X_{n+1} = X_n - \Delta t \nabla \mathcal{F}(X_n) + B e_n.
\]
When the sampling period is sufficiently small the discrete iteration approximates the continuous flow, and the quantization disturbance \(e_n\) plays the role of a stochastic fluctuation.

The equivalence becomes exact when the following identification is made:
\[
F(X,U) \equiv \mathcal{F}(X) = S(X) + R(X) - \langle X, U \rangle.
\]
The entropy potential corresponds to the negative log–likelihood of prediction errors, while the loop–filter potential corresponds to the internal model of the system. The term \(-\langle X,U\rangle\) represents the match between predictions and observations.

Thus, delta–sigma modulation is structurally identical to a variational inference system that performs continuous prediction–error minimization. The quantizer introduces symbolic collapse analogous to discrete perceptual states in cognitive systems. The stability of the modulator corresponds to the existence of a coherent free–energy landscape. Noise shaping corresponds to suppressing prediction error in specific geometric directions. Overload corresponds to the collapse of the variational free–energy structure due to excessive curvature or entropy injection.

This equivalence provides a bridge between engineering, neuroscience, and physics. It shows that delta–sigma modulators are not merely circuits but instances of a universal class of variational feedback systems.

\section{A.46 Global Measurement Geometry and Derived Homotopy Constraints}

The measurement performed by a delta–sigma modulator is an active process. It involves not merely sampling the input signal but actively shaping the entropy landscape through feedback. The derived–geometric framework provides a global characterization of this measurement process. The modulator implements a morphism between derived stacks:
\[
f : \mathcal{X} \longrightarrow \mathcal{Y},
\]
where \(\mathcal{X}\) is the derived signal space and \(\mathcal{Y}\) is the derived symbolic space. The mapping must be homotopically coherent, meaning that the tangent complexes must satisfy compatibility conditions at all stages of the mapping.

This coherence is enforced by the quantizer, which truncates the tangent complex. The loop filter introduces curvature that governs the geometric evolution of the state. Stability requires that the curvature remain within a range compatible with the truncation. When curvature becomes too large, the homotopy condition fails because the projection map discards necessary information. When curvature becomes too small, the projection map cannot enforce coherence and the system drifts.

These constraints define the global measurement geometry of the delta–sigma modulator. The modulator is a derived measurement system in which the symbolic output is the result of a homotopically coherent mapping of continuous geometric states into discrete symbolic states. The entire theory of stability and noise shaping can therefore be understood as the enforcement of homotopy constraints across this mapping.

The next section will develop the general entropy–curvature bridge that unifies all prior results and enables a complete mathematical classification of delta–sigma modulators in geometric–variational terms.

\section{A.47 The Entropy–Curvature Bridge and Universal Delta–Sigma Structure}

The preceding sections have developed three parallel descriptions of delta–sigma modulators: a geometric description based on curvature and entropy flows, a variational description based on free–energy minimization, and a derived–homotopical description based on truncation of tangent complexes. Although each viewpoint illuminates a different aspect of the modulator’s behavior, they are not independent. They are facets of a single underlying principle, which may be called the entropy–curvature bridge. This bridge relates the dissipation of quantization-induced entropy to the curvature properties of the manifold on which the modulator state evolves. Through this relation, all classical and geometric properties of delta–sigma modulation become unified.

Let the entropy potential be denoted by \(S(X)\) and the Riemannian metric by \(g_{ij}(X) = \partial_i \partial_j S(X)\). The curvature tensor \(R_{ijkl}\) derived from this metric encodes how local perturbations of the state grow or contract. The entropy–curvature bridge can be expressed through the identity
\[
\mathrm{div}(\nabla S) = \operatorname{tr}(H_S) = \Delta_g S,
\]
where \(\Delta_g\) is the Laplace–Beltrami operator associated with the entropy metric. This identity shows that the curvature properties of the manifold determine the rate at which entropy dissipates. High positive curvature induces rapid dissipation, while low or negative curvature induces slow dissipation or even amplification of entropy. The balance between entropy injection from quantization and entropy dissipation from curvature therefore governs the stability of the modulator.

The universal structure of delta–sigma modulators emerges from this relationship. The modulator operates by redirecting quantization-induced entropy along geometrically favorable directions, where it can be dissipated by curvature. This redirection is implemented by the loop filter and encoded in the geometry of the free–energy functional. Meanwhile, symbolic collapse induced by quantization ensures that the tangent complex is truncated so that the state trajectory remains within the stable region of the manifold.

Thus, the entropy–curvature bridge unifies the derived, geometric, and variational descriptions, revealing that delta–sigma modulators are governed by an intrinsic principle: stability and noise shaping are the result of curvature–controlled entropy transport.

\section{A.48 Complete Derived Variational Classification Theorem}

The geometric, variational, and homotopical structures of delta–sigma modulators allow the formulation of a classification theorem that encompasses all possible architectures, independent of implementation details. This theorem characterizes a delta–sigma modulator as a triple \((\mathcal{M},\mathcal{F},Q)\), where \(\mathcal{M}\) is an entropy manifold, \(\mathcal{F}\) is a free–energy functional, and \(Q\) is a derived projection. All classical modulators, regardless of order, quantizer structure, or loop-filter configuration, correspond to such a triple.

\medskip

\textbf{Theorem A.6 (Derived Variational Classification).} 
\emph{Let \((\mathcal{M},g)\) be a Riemannian manifold defined by an entropy potential \(S(X)\). Let \(\mathcal{F}(X)\) be a convex free–energy functional of the form}
\[
\mathcal{F}(X) = S(X) + R(X) - \langle X, U \rangle,
\]
\emph{and let \(Q : \mathcal{M} \to \Lambda\) be a projection onto a discrete derived stack \(\Lambda\). Let the state evolve according to}
\[
X_{n+1} = X_n - \Delta t \nabla \mathcal{F}(X_n) + B e_n.
\]
\emph{Then the system is a delta–sigma modulator if and only if:}

(1) \emph{the curvature tensor \(R_{ijkl}\) of \(g\) satisfies the bounded curvature condition},  
(2) \emph{the projection \(Q\) truncates the derived tangent complex in a manner compatible with curvature},  
(3) \emph{the entropy injection from quantization satisfies the bounded truncation condition},  
(4) \emph{the free–energy functional \(\mathcal{F}\) is convex in the stable region and its descent dominates quantization injection},  
(5) \emph{the discrete dynamics admit a compact global attractor invariant under the truncated flow.}

\medskip

\textbf{Proof.}
Necessity follows from the fact that every classical delta–sigma modulator can be expressed in variational form and its stability depends on convexity and curvature conditions. Sufficiency follows by constructing a loop filter that realizes the variational gradient, interpreting the quantizer as the derived projection, and invoking the curvature–entropy balance derived earlier. The global attractor arises from contraction properties of the free–energy functional and bounded quantization injection. \(\square\)

\medskip

This theorem provides the first architecture–independent characterization of delta–sigma modulators. Any feedback system satisfying these five geometric conditions behaves as a delta–sigma modulator, and conversely any delta–sigma modulator satisfies them. This result unifies the field and shows that delta–sigma behavior is a manifestation of deeper geometric and variational laws.

\section{A.49 Geometric Limits of Precision, Resolution, and Symbolic Output}

The resolution of a delta–sigma modulator—the degree to which its bitstream captures the underlying continuous signal—is fundamentally limited by geometric properties of the entropy manifold. Classical descriptions relate resolution to oversampling ratio and loop-filter order. The derived–geometric framework reveals that resolution is determined by the curvature–entropy product introduced earlier.

Let the resolution \(R\) be defined as the effective number of bits in the reconstructed signal. Classical results state that
\[
R \approx L \log_2(\mathrm{OSR}) - C
\]
for some architecture–dependent constant \(C\). The geometric interpretation replaces these analytic quantities with intrinsic geometric invariants. The oversampling ratio corresponds to the temporal resolution of the derived flow, while the order corresponds to the curvature of the manifold near low-frequency modes.

In geometric terms, resolution is governed by the quantity
\[
\Psi = \int_{\mathcal{A}} \frac{1}{K_{\mathrm{min}}(X)} \, d\mu(X),
\]
where \(K_{\mathrm{min}}(X)\) is the minimum sectional curvature at \(X\) and \(\mu\) is the invariant measure on the attractor. The smaller the inverse curvature integrated over the attractor, the higher the resolution. This geometric quantity replaces the classical NTF–based bound.

The quantizer projection affects resolution through symbolic collapse. The projection onto \(\Lambda\) imposes a discrete structure whose granularity imposes a lower bound on the resolution. The derived–geometric theory expresses this limit as a bound on the truncation error of the tangent complex:
\[
\|X - Q(X)\| \ge \delta_{\mathrm{min}},
\]
where \(\delta_{\mathrm{min}}\) is determined by the lattice spacing in the quantizer image. The resolution cannot exceed the inverse of this spacing unless additional geometric contraction is introduced.

Thus resolution is constrained by three geometric factors: curvature, entropy dissipation, and symbolic spacing. These factors define the ultimate limits of precision in delta–sigma converters and reveal that resolution is governed by geometry rather than merely by sampling theory.

\section{A.50 Toward a Final Unified Stability Theorem}

The convergence of geometric, variational, and homotopical principles in the analysis of delta–sigma modulators suggests the existence of a single overarching theorem that encapsulates all known stability results. Such a theorem must unify the entropy–curvature balance, the boundedness of the free–energy functional, the compatibility of the quantizer truncation, and the existence of a global attractor.

The path to this theorem begins with the observation that the stability of a delta–sigma modulator corresponds to the boundedness of the variational evolution of its state. Let the discrete energy difference be written as
\[
\Delta \mathcal{F}_n = \mathcal{F}(X_{n+1}) - \mathcal{F}(X_n).
\]
Stability requires that the expected value of \(\Delta \mathcal{F}_n\) be negative for sufficiently large \(\mathcal{F}(X_n)\). This requirement can be expressed as
\[
\mathbb{E}\left[ \Delta \mathcal{F}_n \right] \le -\alpha \mathcal{F}(X_n) + \beta,
\]
with \(\alpha > 0\). This is the same inequality that arose earlier in the derived Lyapunov context. It implies the existence of a compact sublevel set
\[
\mathcal{K} = \{ X : \mathcal{F}(X) \le C \}
\]
that is invariant under the flow. The attractor lies within \(\mathcal{K}\), and stability corresponds to the persistence of this invariant region.

The final unified stability theorem must establish that the derived tangent truncation, curvature–entropy dissipation, and variational descent jointly guarantee the existence of such a compact invariant region. The next section will express and prove this theorem in full generality, showing that delta–sigma modulators are stable if and only if they satisfy the curvature–entropy compatibility condition.

\section{A.51 The Final Unified Stability Theorem}

The convergence of all previous arguments now allows the formulation of a single theorem that unifies classical stability criteria, Lyapunov theory, curvature–driven geometric dissipation, derived tangent truncation, and variational free–energy descent into one structure. This theorem expresses the necessary and sufficient conditions under which a delta–sigma modulator remains globally stable, regardless of its order, quantizer resolution, or filter architecture. It represents the culmination of the geometric and derived–variational interpretation developed throughout this appendix.

\medskip

\textbf{Theorem A.7 (Final Unified Stability Theorem).}  
\emph{Let \(X_{n+1} = X_n - \Delta t \nabla \mathcal{F}(X_n) + B e_n\) be the state recursion of a delta–sigma modulator on a Riemannian entropy manifold \((\mathcal{M},g)\). Let the free–energy functional be given by \(\mathcal{F}(X) = S(X) + R(X) - \langle X, U\rangle\), where \(S\) is strictly convex and induces the metric \(g\). Let the quantizer projection \(Q\) truncate the derived tangent complex compatibly with curvature. Then the modulator is globally stable if and only if the curvature–entropy compatibility condition}
\[
\langle \nabla S(X), v(X) \rangle \ge \Delta S_Q(X) - \alpha S(X)
\]
\emph{holds for all sufficiently large \(\|X\|\), where \(v\) is the loop–filter vector field, \(\Delta S_Q\) is the entropy injection of quantization, and \(\alpha>0\) is a constant.}

\medskip

\textbf{Interpretation.}  
The pairing \(\langle \nabla S, v\rangle\) measures the rate at which the loop filter dissipates entropy introduced by quantization. The quantity \(\Delta S_Q\) measures the entropy injected by the quantizer. The inequality above states that the dissipative curvature of the manifold must dominate the entropy injection asymptotically. If dissipation dominates injection, the state is pulled inward, producing a compact attractor. If injection dominates dissipation, the state escapes to infinity. This single inequality therefore contains all known stability criteria.

\medskip

\textbf{Proof.}

\emph{Necessity.}  
Assume the modulator is stable. Then the state must remain bounded for all bounded inputs. Let \(X_n\) lie outside a sufficiently large compact region. In this region the dynamics are governed by the competition between the geometric dissipation \(-\nabla S\) and the loop–filter flow \(v\), perturbed by quantization impulses. Because the state remains bounded, the total entropy must not diverge. Thus the expected increment of the entropy potential satisfies
\[
\mathbb{E}\left[S(X_{n+1}) - S(X_n)\right] \le 0.
\]
To first order, expanding in \(\Delta t\),
\[
S(X_{n+1}) - S(X_n) = -\Delta t \langle \nabla S(X_n), \nabla\mathcal{F}(X_n) \rangle + \Delta S_Q(X_n) + O(\Delta t^2).
\]
Since \(\nabla\mathcal{F} = \nabla S + \nabla R - U\), the dominant term in the asymptotic region is \(-\langle\nabla S,\nabla S\rangle\), which is strictly negative. Boundedness therefore requires that the positive contribution from quantization not exceed the negative curvature–induced dissipation. This yields exactly the stated inequality for sufficiently large \(\|X\|\).

\emph{Sufficiency.}  
Assume the curvature–entropy compatibility condition holds. Then for sufficiently large \(\|X\|\),
\[
S(X_{n+1}) - S(X_n) \le -\alpha S(X_n) + \varepsilon
\]
for some small constant \(\varepsilon\) bounding quantization injection. This is precisely the derived Lyapunov inequality developed earlier. Its satisfaction implies that all trajectories eventually enter a compact sublevel set of \(S\) and remain in it indefinitely. This proves global stability.  \(\square\)

\medskip

This theorem formalizes the core geometric principle underlying delta–sigma modulation: stability emerges from a competition between curvature–driven dissipation and entropy injection from quantization. The entire modulator architecture is merely an implementation strategy for shaping this competition so that the dissipative term dominates.

\section{A.52 Curvature–Entropy Compatibility: Detailed Global Proof}

The condition appearing in the unified theorem requires elaboration in order to reveal its geometric significance fully. Let the entropy potential be written in local coordinates as
\[
S(X) = \frac{1}{2} X^\top P(X) X,
\]
where \(P(X)\) is a positive–definite matrix that varies smoothly across the manifold. The gradient is then
\[
\nabla S(X) = P(X) X + \frac{1}{2}\nabla(P(X)) X\otimes X.
\]
The curvature tensor \(R_{ijkl}\) associated with this metric determines how \(\nabla S\) evolves as the state moves. In particular, directional curvature along \(v(X)\) is given by
\[
K_v(X) = \frac{R(\nabla S, v, \nabla S, v)}{\|\nabla S \wedge v\|^2}.
\]

Entropy dissipation along the vector field is expressed as 
\[
D(X) = \langle\nabla S(X), v(X)\rangle.
\]
This quantity measures how rapidly the loop filter drives the state into regions of lower entropy. The quantization injects entropy in discrete impulses, and the expected injection per step is \(\Delta S_Q(X)\). The curvature–entropy compatibility condition states that
\[
D(X) \ge \Delta S_Q(X) - \alpha S(X)
\]
for sufficiently large \(\|X\|\). The proof divides into two complementary regimes.

When \(X\) lies in a region of high curvature, the quantity \(D(X)\) grows superlinearly with \(\|X\|\), while the quantization injection grows only linearly. Thus dissipation dominates injection, and the inequality holds naturally. This explains the stability of high–order modulators at moderate amplitudes.

When \(X\) lies in a region of low curvature, \(D(X)\) may decrease or even vanish. In this regime the \(-\alpha S(X)\) term becomes essential. Strict convexity of the entropy potential ensures that \(\alpha S(X)\) grows without bound as \(\|X\|\to\infty\). Therefore, even if \(D(X)\) is small, the combined term still dominates the bounded quantization injection. This ensures boundedness even when the vector field becomes weak.

Thus the compatibility condition reflects a dual guarantee: curvature dominates quantization in high–curvature regions, while the convexity of the entropy potential dominates quantization in low–curvature regions. The interplay of these two mechanisms produces the global attractor.

\section{A.53 Consequences for Architecture Design}

The unified stability theorem and curvature–entropy compatibility condition have immediate implications for architecture design. They reveal that stability is not a property of specific loop–filter coefficients or quantizer resolutions but a geometric property of the free–energy manifold. The entire modulator architecture must therefore be understood as a mechanism for shaping curvature so that dissipation uniformly dominates injection.

One consequence is that placing too many zeros at \(z=1\) introduces excessive curvature near low–frequency modes. While this improves noise–shaping performance, it can violate the compatibility condition in high–amplitude regions, causing overload. The derived theory shows that the classical “order limit” arises not from algebraic complications but from the breakdown of curvature–entropy compatibility.

Another consequence is that multi–bit quantization improves stability not because it reduces quantization noise in the conventional sense, but because it reduces \(\Delta S_Q\), lowering the entropy injection term in the compatibility inequality. Thus multi–bit modulators enlarge the region of compatibility and allow steeper curvature without risking instability.

A third consequence is that dithering improves stability by reducing the worst–case symbolic distortion produced by quantization collapse. This smooths the pushforward of the attractor measure and reduces peaks in entropy injection. From the derived viewpoint, dithering enriches the tangent complex so that collapse becomes gentler, making the compatibility condition easier to satisfy.

These design implications are direct corollaries of the unified theorem and illustrate the power of the geometric formulation.

\section{A.54 Extensions to Continuous–Time and MIMO Systems}

The curvature–entropy compatibility condition extends naturally to continuous–time modulators. Let the continuous state evolve as
\[
\frac{dX}{dt} = -\nabla\mathcal{F}(X) + \eta(t),
\]
with quantization impulses modeled as a distributional term \(\eta(t)\). The energy balance equation becomes
\[
\frac{d}{dt} S(X(t)) = -\langle\nabla S(X), \nabla\mathcal{F}(X)\rangle + \Xi(t),
\]
where \(\Xi(t)\) is the instantaneous entropy injection. The continuous–time version of the compatibility condition is
\[
\langle\nabla S(X), \nabla\mathcal{F}(X)\rangle \ge \Xi(t) - \alpha S(X),
\]
identical in form to the discrete version. Continuous–time modulators therefore obey the same geometric laws.

For MIMO modulators the state manifold becomes a product manifold \(\mathcal{M}^{(1)} \times \cdots \times \mathcal{M}^{(m)}\), with a block–diagonal or fully coupled entropy metric. Stability requires that the curvature–entropy compatibility condition hold on each manifold component and on all mixed curvature terms. Cross–coupled entropy injection between channels must be dissipated by curvature in orthogonal directions. This yields a multidimensional version of the unified stability theorem.

Thus the same geometric and variational laws hold for continuous–time and MIMO modulators, demonstrating the universality of the entropy–curvature bridge.

\section{A.55 Closing Remarks and Integration with RSVP}

The geometric framework developed in this appendix reveals that delta–sigma modulators belong to a far broader category of systems governed by entropy transport on curved manifolds. In the RSVP theoretical framework entropy flows, vector fields, and symbolic collapse form the basis of physical and cognitive dynamics. Delta–sigma modulators manifest the same principles at a smaller scale. Their stability emerges from the same balance of scalar entropy descent, vector–field flow, and discrete symbolic collapse that governs the RSVP dynamics of perception, cognition, and measurement.

Thus delta–sigma modulation is more than an engineering construction. It is a canonical instance of entropy–shaping dynamics on a curved manifold, providing a controlled and mathematically transparent microcosm of the RSVP field equations. The geometric theory developed here will therefore reappear in the later chapters of this book, where delta–sigma modulation becomes a template for measurement, inference, and coherence in natural and artificial systems.

